\section{Security}\label{sec:security}
%----------------------------------------------------------------------

We prove that the zkVM construction is \emph{correct-trace extractable}: any
adversary that produces a convincing final proof can be used to extract an
explicit, valid execution trace.  The proof proceeds by descending the recursion
tree layer by layer, reducing everything ultimately to the knowledge soundness
of the inner-circuit SNARKs, the binding properties
of the memory commitment scheme $\Com_{\mathsf{mem}}$, and the collision-resistance of $\Com_{\mathsf{bus}}$.

\begin{remark}[Idealized proof model]\label{rem:idealized}
Our proof is idealized/heuristic in two respects.

First, Definition~\ref{def:extractable} below postulates a straight-line
extractor that runs in essentially the time of the adversary, with only
additive $\mathsf{poly}(\lambda)$ overhead. Because our SNARKs are hash-based,
the natural model is the random oracle model (ROM), in which such
straight-line extractors do exist. The subtlety lies in the recursion:
composing straight-line extraction across the layers of the proof tree (as in
proof-carrying data, PCD) requires each SNARK to remain knowledge-sound
\emph{relative to the same random oracle}, i.e.\ to be a \emph{relativized}
argument system, and such relativized SNARKs provably do not exist in general
\cite{EPRINT:relativized}. Our analysis idealizes over this obstruction: we assume that straight-line extraction composes across the
recursion, effectively putting the non-existence of relativized SNARKs under
the rug.

Second, while the analysis does produce explicit advantage bounds and reduction
running times (Theorem~\ref{thm:main}), the idealization above means these
should not be read as a concrete security level: they are meaningful as a
statement about the reduction structure --- the recursion topology is sound and
the circuit relations compose consistently --- but a concrete security claim
would require replacing the idealized straight-line extraction assumption with
an analysis of the actual instantiation.
\end{remark}

%----------------------------------------------------------------------
\subsection{Definitions}\label{sec:security:defs}
%----------------------------------------------------------------------

We work with knowledge-sound argument systems admitting straight-line
extraction.  Let $\lambda$ denote the security parameter;
``PPT'' means probabilistic polynomial time in $\lambda$.

\begin{definition}[Knowledge-soundness]\label{def:extractable}
A non-interactive argument system $\Pi=(\mathsf{Prove},\mathsf{Verify})$ for
relation $\mathcal{R}$ is \emph{knowledge-sound} if there exists a PPT
algorithm $\mathcal{E}$ (the \emph{extractor}) such that for every PPT
adversary $\mathcal{A}$:
\begin{equation}\label{eq:extractable}
  \mathsf{Adv}^{\mathsf{ks}}_{\Pi}(\mathcal{A}) \;:=\;
  \Pr\!\Bigl[
    (x,\pi)\leftarrow\mathcal{A}(1^\lambda);\; w\leftarrow\mathcal{E}(x,\pi)
    \;:\; \Pi.\mathsf{Verify}(x,\pi)=1 \;\wedge\; (x;\,w)\notin\mathcal{R}
  \Bigr] \leq \mathsf{negl}(\lambda).
\end{equation}

We assume that the running time of $\mathcal{E}(x,\pi)$ is at most
$\mathsf{Time}(\mathcal{A}(1^\lambda)) + \mathsf{poly}(\lambda)$ for every
adversary $\mathcal{A}$ and every $(x,\pi)$ in the support of $\mathcal{A}$'s
output, i.e.\ the
extractor incurs only an additive polynomial overhead over a single evaluation of
$\mathcal{A}$ on the given proof.
We call $\mathsf{Adv}^{\mathsf{ks}}_{\Pi}(\mathcal{A})$ the knowledge-soundness
error of $\mathcal{A}$ against $\Pi$.
\end{definition}

\begin{remark}[Straight-line extraction]
The extractor $\mathcal{E}$ above is straight-line: it takes only the
statement $x$ and the proof $\pi$ as input and does not rewind $\mathcal{A}$. This can actually not be achieved if $\pi$ is succinct, but this is part of our (over-)idealization, see Remark~\ref{rem:idealized} above. Namely, Definition~\ref{def:extractable} is
an idealized requirement that cannot be satisfied by any standard construction
in the plain model (see Remark~\ref{rem:idealized}).  It serves as a
\emph{proxy}: we state it formally so that the reduction structure is
explicit, with the understanding that concrete instantiations will replace it
with ROM-based extractability at the cost of the overhead
discussed in Remark~\ref{rem:idealized}.
\end{remark}

\begin{instruction}
  In the following, we omit commitment keys for simplicity of exposition. Teams should add appropriate definitions and adaptations if they rely on commitment keys.
\end{instruction}

\begin{instruction}
  We also leave out completeness definitions, as the document should just illustrate the structure of a potential soundness proof. Teams should add completeness arguments as well.
\end{instruction}

\begin{definition}[Position-binding and update-binding vector commitment]\label{def:binding}
A vector commitment scheme
$\Com=(\mathsf{Commit},\mathsf{Open},\mathsf{Verify})$
satisfies the two independent properties below --- \emph{position-binding} and
\emph{update-binding} --- which a scheme may have separately or together.
\begin{enumerate}
  \item \textbf{Position-binding.} For every PPT adversary $\mathcal{A}$, the following is negligible:
\[
  \mathsf{Adv}^{\mathsf{pos}}_{\Com}(\mathcal{A}) \;:=\;
  \Pr\!\left[\begin{array}{l}
    (C,i,v,\pi,v',\pi')\leftarrow\mathcal{A}(1^\lambda) \;:\; \\[4pt]
    \mathsf{Verify}(C,i,v,\pi)=1 \;\wedge\;
    \mathsf{Verify}(C,i,v',\pi')=1 \;\wedge\; v\neq v'
  \end{array}\right].
\]
  \item \textbf{Update binding.} For every PPT adversary $\mathcal{A}$, the
    following is negligible, where $C:=\mathsf{Commit}(\vec{m})$:
\[
  \mathsf{Adv}^{\mathsf{upd}}_{\Com}(\mathcal{A}) \;:=\;
  \Pr\!\left[\begin{array}{l}
    (\vec{m}, \addr, x, C', \pi) \leftarrow \mathcal{A}(1^\lambda) \;:\; \\[4pt]
    \mathsf{Verify}(C,\addr,\vec{m}[\addr],\pi)=1 \;\wedge\;
    \mathsf{Verify}(C',\addr,x,\pi)=1 \\[4pt]
    {}\wedge\; C' \neq \mathsf{Commit}(\vec{m}[\addr\mapsto x])
  \end{array}\right].
\]
In other words, a single opening proof $\pi$ that opens an honest commitment
$C=\mathsf{Commit}(\vec{m})$ at position $\addr$ and also opens $C'$ at $\addr$
to $x$ forces $C'$ to be the honest commitment of the updated vector
$\vec{m}[\addr\mapsto x]$. In particular $C'$ is then itself a well-formed
commitment --- an actual output of $\mathsf{Commit}$ --- and not merely a value
that verifies against some proofs.
\end{enumerate}
\end{definition}
\begin{remark}[Merkle trees and update binding]
Standard Merkle tree commitments satisfy both position-binding and update
binding under collision-resistance of the underlying hash function.  The
position-binding property is immediate.  For update binding, a Merkle
authentication path $\pi$ for position $\addr$ consists of the sibling hashes
along the root-to-leaf path; since $C=\mathsf{Commit}(\vec{m})$ is honest, those
siblings must --- except for a hash collision --- be the honest siblings of
$\vec{m}$, so recomputing the root with the new leaf $x$ yields exactly
$\mathsf{Commit}(\vec{m}[\addr\mapsto x])$, and no other value of $C'$ can pass.
We leave out a formal proof for now.
\end{remark}

\begin{definition}[Collision-resistant bus commitment]\label{def:bus-cr}
A bus commitment scheme $\Com_{\mathsf{bus}}$ is \emph{collision-resistant} if for every PPT adversary $\mathcal{A}$, the following is negligible:
\[
\mathsf{Adv}^{\mathsf{cr}}_{\Com_{\mathsf{bus}}}(\mathcal{A}):= \Pr\!\left[\begin{array}{l}
    (\widehat{\BBB},\widehat{\BBB}') \leftarrow \mathcal{A}(1^\lambda) \;:\; \\[4pt]
    \widehat{\BBB}\neq\widehat{\BBB}' \;\wedge\;
    \Com_{\mathsf{bus}}(\widehat{\BBB})=\Com_{\mathsf{bus}}(\widehat{\BBB}')
  \end{array}\right].
\]
\end{definition}



\begin{definition}[zkVM system]\label{def:zkvm}
A \emph{zkVM system} is a tuple
\[
  \mathsf{zkVM} = \bigl(
    \{\varphi_{\op}\}_{\op},\;
    \Com_{\mathsf{bus}},\;
    \Com_{\mathsf{mem}},\;
    \Pi_0,\, \Pi_1,\, \Pi_2,\, \Pi_3,\, \Pi_4
  \bigr)
\]
consisting of the following components.
\begin{enumerate}
  \item \textbf{Step predicates.}  For each operation $\op$ (including
    $\step$, $\rea$, $\writ$, $\keccak$, $\poseidon$, $\range$), a
    predicate $\varphi_{\op}$ over VM states, expressible as a system of
    polynomial constraints over the field (as defined in
    Section~\ref{sec:proof-arch}).

  \item \textbf{Bus commitment scheme.}
    $\Com_{\mathsf{bus}}$ is a collision-resistant
    bus commitment scheme (Definition~\ref{def:bus-cr}) used to commit to the shared bus
    $\mathcal{B}$ in the inner-circuit relations $\relation_{0,j}$.

  \item \textbf{Memory commitment scheme.}
    $\Com_{\mathsf{mem}} = (\mathsf{Commit}, \mathsf{Open}, \mathsf{Verify})$,
    a position-binding and update-binding vector commitment scheme
    (Definition~\ref{def:binding}) used to represent the memory component
    of each committed VM state $\hat{S}$.

  \item \textbf{SNARK hierarchy.}
    \begin{itemize}
      \item $\Pi_0 = \{\Pi_{0,j}\}_{j \in \{\step,\keccak,\poseidon,\range\}}$:
        the four \emph{inner-circuit} SNARKs for relations
        $$\relation_{0,\step},\, \relation_{0,\keccak},\,
         \relation_{0,\poseidon},\, \relation_{0,\range},$$ respectively.
      \item $\Pi_1$: the \texttt{segment} SNARK for relation $\relation_1$
        (verifies the four inner proofs).
      \item $\Pi_2$: the \texttt{convert} SNARK for relation $\relation_2$
        (1-to-1 normalization).
      \item $\Pi_3$: the \texttt{combine} SNARK for relation $\relation_3$
        (2-to-1 recursive merging).
      \item $\Pi_4$: the \texttt{embed} SNARK for relation $\relation_4$
        (final proof).
    \end{itemize}
    Each $\Pi_i = (\Pi_i.\mathsf{Prove}, \Pi_i.\mathsf{Verify})$ is a
    non-interactive argument system for the respective relation.
\end{enumerate}
\end{definition}

\begin{definition}[Correct-execution relation]\label{def:relation-star}
Fix a program code $\code$ and a step count $T$ satisfying $N_{\seg}\mid T$,
$T\ge 2N_{\seg}$, and $T\le\mathsf{poly}(\lambda)$ (system parameters, as in the
\texttt{embed} circuit of Section~\ref{sec:proof-arch}); neither $\code$ nor $T$
is chosen by the adversary.

The \emph{correct-execution relation} $\relation^*=\relation^*_{\code,T}$ has as
statements boundary pairs $(S_0,S_T)$ of full-memory VM states
$S=(\pc,\regs,\mem)$, and as witnesses candidate full-memory traces
$(S_0',\ldots,S_T')$:
\begin{equation}\label{eq:relation-star}
  \relation^* \;=\; \left\{\,\bigl((S_0,S_T);\;(S_0',\ldots,S_T')\bigr)\;\middle|\;
    S_0'=S_0 \;\wedge\; S_T'=S_T \;\wedge\;
    \forall i\in[T]:\varphi_{\step}(S_{i-1}',S_i')=\mathrm{TRUE}
  \right\},
\end{equation}
where $\varphi_{\step}$ is the step predicate of Section~\ref{sec:proof-arch}
(the disjunction over all operation predicates $\varphi_{\op}$).
\end{definition}

\begin{definition}[Final argument system]\label{def:final-as}
The \emph{final argument system} of $\mathsf{zkVM}$ is the non-interactive
argument system $\Pi^*=(\Pi^*.\mathsf{Prove},\Pi^*.\mathsf{Verify})$ for the
correct-execution relation $\relation^*$ (Definition~\ref{def:relation-star})
whose proofs are the final \texttt{embed} proofs $\pi^4$ and whose verification
algorithm, on statement $(S_0,S_T)$ and proof $\pi^4$, proceeds as follows:
\begin{enumerate}
  \item set
    $$\hat{S}_0 := (\pc_0,\regs_0,\mathsf{Commit}(\mem_0)),\quad
    \hat{S}_T := (\pc_T,\regs_T,\mathsf{Commit}(\mem_T));$$
  \item output $\Pi_4.\mathsf{Verify}((\hat{S}_0,\hat{S}_T),\pi^4)$,
\end{enumerate}
where $\mathsf{Commit}$ is that of $\Com_{\mathsf{mem}}$.  Its prover
$\Pi^*.\mathsf{Prove}$ is the honest proving pipeline of
Section~\ref{sec:proof-arch} instantiated at $(\code,T)$; it plays no role in
what follows.\footnote{Observe that the definition of knowledge soundness just
depends on a verification algorithm and not on the associated prover.}
\end{definition}

\begin{definition}[Correct-trace extractability]\label{def:cte}
The zkVM $\mathsf{zkVM}$ is \emph{correct-trace extractable} if its final
argument system $\Pi^*$ (Definition~\ref{def:final-as}) is knowledge-sound for
the correct-execution relation $\relation^*$ (Definition~\ref{def:relation-star})
in the sense of Definition~\ref{def:extractable}.  We write
\[
  \mathsf{Adv}^{\mathsf{cte}}_{\mathsf{zkVM}}(\mathcal{A})
  \;:=\;\mathsf{Adv}^{\mathsf{ks}}_{\Pi^*}(\mathcal{A})
\]
for the \emph{CTE advantage} of a PPT adversary $\mathcal{A}$, and
$\mathcal{E}_{\mathsf{zkVM}}$ for the extractor whose existence is asserted; the
additive running-time convention of Definition~\ref{def:extractable} applies to
it unchanged.
\end{definition}

\begin{remark}[Unrolling Definition~\ref{def:cte}]\label{rem:cte-experiment}
Written out, Definition~\ref{def:cte} says that there is a PPT extractor
$\mathcal{E}_{\mathsf{zkVM}}$ such that for every PPT adversary $\mathcal{A}$ the
experiment
\[
  \mathsf{Exp}_{\mathsf{zkVM},\mathcal{E}_{\mathsf{zkVM}},\mathcal{A}}^{\mathsf{cte}}(\lambda)
\]
below returns $1$ with at most negligible probability.
\begin{enumerate}
  \item \textbf{Adversary.} $\mathcal{A}(1^\lambda)$ outputs a statement
    $(S_0,S_T)$ --- an initial state $S_0=(\pc_0,\regs_0,\mem_0)$ and a final
    state $S_T=(\pc_T,\regs_T,\mem_T)$ --- together with a proof $\pi^4$; the
    program code $\code$ and step count $T$ are fixed system parameters.
  \item \textbf{Verification.} If $\Pi^*.\mathsf{Verify}((S_0,S_T),\pi^4)=0$
    (Definition~\ref{def:final-as}), return $0$.
  \item \textbf{Extraction.} Run $\mathcal{E}_{\mathsf{zkVM}}(S_0,S_T,\pi^4)$ to
    obtain full-memory VM states $S_0',S_1',\ldots,S_T'$, where
    $S_i'=(\pc_i',\regs_i',\mem_i')$.
  \item \textbf{Winning condition.} Return $1$ (adversary wins) if and only if
    $\bigl((S_0,S_T);(S_0',\ldots,S_T')\bigr)\notin\relation^*$, i.e.\ if any of
    the following hold:
    \begin{itemize}
      \item $\pc_0' \neq \pc_0$, $\regs_0' \neq \regs_0$, or
        $\mem_0' \neq \mem_0$; or
      \item $\pc_T' \neq \pc_T$, $\regs_T' \neq \regs_T$, or
        $\mem_T' \neq \mem_T$; or
      \item $\varphi_{\step}(S_i', S_{i+1}') \neq \mathrm{TRUE}$ for some
        $i \in \{0,\ldots,T-1\}$.
    \end{itemize}
\end{enumerate}
This is Definition~\ref{def:extractable} instantiated at $(\relation^*,\Pi^*)$
verbatim; in particular the boundary checks are not extra requirements bolted
onto the game, but precisely the two boundary conjuncts
of~\eqref{eq:relation-star}.  Note also that a statement of $\relation^*$
carries the full memory vectors $\mem_0,\mem_T$, so $\Pi^*.\mathsf{Verify}$ is
not succinct; succinctness is a property of $\Pi_4$ and is orthogonal to
correct-trace extractability.
\end{remark}

%----------------------------------------------------------------------
\subsection{Memory extractability}\label{sec:security:bus}
%----------------------------------------------------------------------
The inner circuits operate on committed-memory states $\hat{S}$ and a bus
$\widehat{\BBB}$, whereas the correct-execution relation
$\relation^*$~\eqref{eq:relation-star} requires a valid
execution trace in terms of full-memory states $S$ with \emph{explicitly
reconstructed} memory vectors $\mem_k$.  The following proposition bridges
this gap and shows that, once the commitment openings are embedded in the
SNARK witnesses, the bus-based predicate implies the bus-free predicate:
introducing a bus does not weaken the correctness guarantee.  In particular,
position-binding and update-binding of $\Com_{\mathsf{mem}}$ bound the
probability that the extracted full-memory trace deviates from the
committed-memory execution recorded by the SNARKs.

It does so in three respects: (i)~it eliminates the bus; (ii)~it lifts the
step predicate from committed-memory states to full-memory states whose
values at accessed addresses are uniquely determined by the position-binding
commitment openings embedded in the SNARK witness; and (iii)~it carries the
commitment invariant $\hat{S}.\mem=\mathsf{Commit}(\mem)$ forward by one step,
so that the invariant is a \emph{conclusion} of the proposition rather than a
side condition established elsewhere.  Concretely, it shows
that if the bus-parameterised step predicate $\widehat{\varphi}_{\step}$
holds (as guaranteed by Lemma~\ref{lem:segment}), then the bus-free,
full-memory predicate $\varphi_{\step}(S_1,S_2)$ holds for the states
$S_j := (\hat{S}_j.\pc,\hat{S}_j.\regs,\mem_j)$ formed with the reconstructed
memory vectors, and $\hat{S}_2.\mem=\mathsf{Commit}(\mem_2)$; the required
memory values are extracted from the commitment openings embedded in the
definitions of $\widehat{\varphi}_{\rea}$ and $\widehat{\varphi}_{\writ}$.
It is invoked once per step in Step~6 of the proof of
Theorem~\ref{thm:main}, where the two conclusions are used respectively as the
per-step guarantee required by the step conjunct
of~\eqref{eq:relation-star} and as the induction
hypothesis for the next step.

The second memory vector is not a free choice of the adversary: it is obtained
from the first by the deterministic \emph{memory reconstruction rule}
\begin{equation}\label{eq:mem-reconstruct}
  \mem_2 \;:=\;
  \begin{cases}
    \mem_1[\,\hat{S}_1.\regs[0]\mapsto\hat{S}_1.\regs[1]\,]
      & \text{if } \code[\hat{S}_1.\pc]=\writ,\\[2pt]
    \mem_1 & \text{otherwise,}
  \end{cases}
\end{equation}
which is exactly the rule the extractor of Theorem~\ref{thm:main} applies when
it reconstructs the full-memory trace.  Fixing $\mem_2$ this way is what makes
the commitment invariant provable rather than assumable: an adversary that
could choose $\mem_2$ freely subject to $\hat{S}_2.\mem=\mathsf{Commit}(\mem_2)$
would already have been handed the invariant as a hypothesis, and the
update-binding term in the bound below would do no work.

Throughout, we write
\[
  \widehat{\varphi}_{\step}(\hat{S}_i,\hat{S}_{i+1},\widehat{\BBB})
\]
as shorthand for the four-argument predicate~\eqref{eq:step-bus2} instantiated
with the memory-opening witness $\pi^{\mathsf{mem}}_i$ that is carried alongside
the states:
\[
  \widehat{\varphi}_{\step}(\hat{S}_i,\hat{S}_{i+1},\widehat{\BBB},\pi^{\mathsf{mem}}_i)=\mathrm{TRUE}.
\]
The shorthand is \emph{not} an existential quantification over
$\pi^{\mathsf{mem}}_i$: the witness is a fixed entry of the extracted
$\relation_{0,\step}$ witness vector.  Wherever the shorthand and an explicit
memory-opening witness both occur --- as in
Proposition~\ref{prop:memory-extractability}, whose second and fourth
hypotheses each refer to $\pi^{\mathsf{mem}}$ --- they denote the same object,
so the disjunct of $\widehat{\varphi}_{\step,\bus}$ that is satisfied is the one
evaluated on that witness.

\begin{proposition}[Memory extractability]
  \label{prop:memory-extractability}
For every PPT adversary $\mathcal{A}$ that outputs a tuple
$$(\hat{S}_1,\hat{S}_2,\mem_1,\widehat{\BBB},\pi^{\mathsf{mem}})$$
such that, defining $\mem_2$ by the reconstruction
rule~\eqref{eq:mem-reconstruct} and writing
$S_j:=(\hat{S}_j.\pc,\hat{S}_j.\regs,\mem_j)$, the following hold with
probability $\varepsilon$:
\begin{itemize}
  \item $\hat{S}_1.\mem = \mathsf{Commit}(\mem_1)$;
  \item $\widehat{\varphi}_{\step}(\hat{S}_1,\hat{S}_2,\widehat{\BBB})=\mathrm{TRUE}$;
  \item the step fails to transfer, i.e.\ the commitment invariant is not
    carried forward or the full-memory step predicate does not hold:
    \[
      \hat{S}_2.\mem \neq \mathsf{Commit}(\mem_2)
      \qquad\text{or}\qquad
      \varphi_{\step}(S_1,S_2)\neq\mathrm{TRUE};
    \]
  \item the memory-opening witness $\pi^{\mathsf{mem}}$ satisfies,
    depending on the operation type:
    \begin{itemize}
      \item \emph{Read}: $\pi^{\mathsf{mem}} = \pi$ with
        $\hat{S}_1.\mem=\hat{S}_2.\mem$ and
        $\mathsf{Verify}(\hat{S}_1.\mem,\hat{S}_1.\regs[0],\hat{S}_2.\regs[1],\pi)=1$;
      \item \emph{Write}: $\pi^{\mathsf{mem}} = (\pi_{\writ}, v_{\mathsf{old}})$ with
        \begin{align*}
        &\mathsf{Verify}(\hat{S}_1.\mem,\hat{S}_1.\regs[0],v_{\mathsf{old}},\pi_{\writ})=1,\\
        &\mathsf{Verify}(\hat{S}_2.\mem,\hat{S}_1.\regs[0],\hat{S}_1.\regs[1],\pi_{\writ})=1;
        \end{align*}
      \item \emph{Other}: no condition is imposed on $\pi^{\mathsf{mem}}$;
    \end{itemize}
\end{itemize}
there exist PPT reduction adversaries $\mathcal{D}^{\mathsf{pos}}$ and
$\mathcal{D}^{\mathsf{upd}}$, each constructed from $\mathcal{A}$ and running
in time $\mathsf{Time}(\mathcal{A})+\mathsf{poly}(\lambda)$, such that
\[
  \varepsilon \leq\;
  \mathsf{Adv}^{\mathsf{pos}}_{\Com_{\mathsf{mem}}}(\mathcal{D}^{\mathsf{pos}})
  + \mathsf{Adv}^{\mathsf{upd}}_{\Com_{\mathsf{mem}}}(\mathcal{D}^{\mathsf{upd}}).
\]
Equivalently: whenever the first, second and fourth hypotheses hold, both
$\hat{S}_2.\mem=\mathsf{Commit}(\mem_2)$ and
$\varphi_{\step}(S_1,S_2)=\mathrm{TRUE}$, except with probability at most
$\mathsf{Adv}^{\mathsf{pos}}_{\Com_{\mathsf{mem}}}(\mathcal{D}^{\mathsf{pos}})
+\mathsf{Adv}^{\mathsf{upd}}_{\Com_{\mathsf{mem}}}(\mathcal{D}^{\mathsf{upd}})$.
\end{proposition}

\begin{remark}[On the fourth hypothesis]\label{rem:mem-ext-hyp4}
Two points about the case distinction in the fourth hypothesis.

First, the \emph{Other} case is deliberately unconstrained.  The remaining
disjuncts of~\eqref{eq:step-bus2} never read $\pi^{\mathsf{mem}}$, so nothing in
the construction forces its value on a non-memory step, and the extractor of
Lemma~\ref{lem:segment} returns whatever the prover placed in that witness
field.  Demanding $\pi^{\mathsf{mem}}=\bot$ there would be a requirement the
zkVM does not enforce, and would make the proposition inapplicable to exactly
those steps; the proof never uses it.

Second, which of the three cases applies is determined by
$\op:=\code[\hat{S}_1.\pc]$ and is therefore not a free choice of
$\mathcal{A}$.  Every $\varphi'_{\op}$ carries the fetch conjunct
$\code[\pc_1]=\op$ of~\eqref{eq:phiop} --- for the two precompiles it is carried
by the chip predicate $\widehat{\varphi}_{\op}(\widehat{\BBB})$ rather than by
the bus-membership disjunct itself --- so under the second hypothesis at most
one disjunct of $\widehat{\varphi}_{\step,\bus}$ can be satisfied.  The
operation type is thus well defined, and the case split in the proof below is
the same case split as the one above.
\end{remark}

\begin{proof}
We define two bad events, each a violation of one binding property of
$\Com_{\mathsf{mem}}$, and bound each by a single reduction that runs
$\mathcal{A}$ once. We then show (Step~A) that off both events the assumption
$\widehat{\varphi}_{\step}(\hat{S}_1,\hat{S}_2,\widehat{\BBB})=\mathrm{TRUE}$
forces both $\hat{S}_2.\mem=\mathsf{Commit}(\mem_2)$ and
$\varphi_{\op}(S_1,S_2)=\mathrm{TRUE}$ for the operation $\op$ selected by
the satisfied disjunct, and hence $\varphi_{\step}(S_1,S_2)=\mathrm{TRUE}$ ---
contradicting the third winning condition. A union bound then gives the claim.

\paragraph{Bad events.}
Run $\mathcal{A}(1^\lambda)$ to obtain
$(\hat{S}_1,\hat{S}_2,\mem_1,\widehat{\BBB},\pi^{\mathsf{mem}})$ and set $\mem_2$
by~\eqref{eq:mem-reconstruct}. By the first hypothesis
$\hat{S}_1.\mem=\mathsf{Commit}(\mem_1)$, correctness of
$\Com_{\mathsf{mem}}$ makes the honest opening $\mathsf{Open}(\mem_1,p)$ verify
against $\hat{S}_1.\mem$ at every position $p$.

All memory vectors have the same length $n$, fixed by the system parameters,
and positions range over the index domain $[n]$.  Writing
$\addr:=\hat{S}_1.\regs[0]$ and $x:=\hat{S}_1.\regs[1]$, let
\[
  \mathcal{V}\;:=\;\{\,\mem_1,\;\mem_2\,\}
\]
be the \emph{candidate vectors}.  Both have length $n$ and are computable in
time $\mathsf{poly}(\lambda)$ from $\mathcal{A}$'s output --- $\mem_2$ because
the reconstruction rule~\eqref{eq:mem-reconstruct} is deterministic and
explicit --- so a reduction may open either of them honestly at any position.
Note that on a write step $\mem_2$ \emph{is} the updated vector
$\mem_1[\addr\mapsto x]$, so $\mathcal{V}$ already contains every vector the
argument below opens.  The well-formedness conjunct of $\varphi'_{\rea}$ and
$\varphi'_{\writ}$ guarantees $\addr\in[n]$ on the two memory operations, so
$\mem_1[\addr]$, $\mathsf{Open}(\mem_1,\addr)$ and $\mem_1[\addr\mapsto x]$ are
all well defined there.  Define:
\begin{itemize}
  \item $E_{\mathsf{pos}}$ (\emph{position-binding collision}): there are a root
    $R\in\{\hat{S}_1.\mem,\hat{S}_2.\mem\}$, a position $p$, and two accepting
    openings of $R$ at $p$ to distinct values, each either read from
    $\pi^{\mathsf{mem}}$ or computed as an honest opening $\mathsf{Open}(\vec{v},p)$
    of a candidate vector $\vec{v}\in\mathcal{V}$.
  \item $E_{\mathsf{upd}}$ (\emph{update-binding collision}):
    $\pi^{\mathsf{mem}}=(\pi_{\writ},v_{\mathsf{old}})$ is a write witness whose
    shared path $\pi_{\writ}$ opens the honest commitment $\hat{S}_1.\mem$ at
    $\addr:=\hat{S}_1.\regs[0]$ to $\mem_1[\addr]$ and opens $\hat{S}_2.\mem$ at
    $\addr$ to $x:=\hat{S}_1.\regs[1]$, yet
    $\hat{S}_2.\mem\neq\mathsf{Commit}(\mem_1[\addr\mapsto x])$.
\end{itemize}

\paragraph{The two reductions.}
$\mathcal{D}^{\mathsf{pos}}$ runs $\mathcal{A}(1^\lambda)$ once, then compares the
candidate vectors in $\mathcal{V}$ pairwise and against the values opened by
$\pi^{\mathsf{mem}}$; whenever
$E_{\mathsf{pos}}$ occurs it computes the corresponding honest openings and
outputs the two openings
$(R,p,v,\rho)$ and $(R,p,v',\rho')$ with $v\neq v'$ to the position-binding
challenger, breaking it. Hence
$\Pr[E_{\mathsf{pos}}]\le\mathsf{Adv}^{\mathsf{pos}}_{\Com_{\mathsf{mem}}}(\mathcal{D}^{\mathsf{pos}})$.
$\mathcal{D}^{\mathsf{upd}}$ runs $\mathcal{A}(1^\lambda)$ once; whenever
$E_{\mathsf{upd}}$ occurs it outputs
\[
  (\mem_1,\addr,x,\hat{S}_2.\mem,\pi_{\writ})
\]
to the update-binding challenger, breaking it: $\hat{S}_1.\mem=\mathsf{Commit}(\mem_1)$
is honest, the shared path $\pi_{\writ}$ opens it at $\addr$ to $\mem_1[\addr]$
and opens $\hat{S}_2.\mem$ at $\addr$ to $x$, yet
$\hat{S}_2.\mem\neq\mathsf{Commit}(\mem_1[\addr\mapsto x])$. Hence
$\Pr[E_{\mathsf{upd}}]\le\mathsf{Adv}^{\mathsf{upd}}_{\Com_{\mathsf{mem}}}(\mathcal{D}^{\mathsf{upd}})$.
Each reduction makes one run of $\mathcal{A}$, a linear scan of the candidate
vectors and a constant number of honest openings, so both run in time
$\mathsf{Time}(\mathcal{A})+\mathsf{poly}(\lambda)$ (note that $\mathcal{A}$
itself outputs $\mem_1$ and hence already runs in time $\Omega(n)$).

\paragraph{Step~A: off both events, the invariant and the step predicate transfer.}
\textit{Goal:} assuming $\neg E_{\mathsf{pos}}\wedge\neg E_{\mathsf{upd}}$,
\[
  \widehat{\varphi}_{\step}(\hat{S}_1,\hat{S}_2,\widehat{\BBB})=\mathrm{TRUE}
\]
implies $\hat{S}_2.\mem=\mathsf{Commit}(\mem_2)$ and
$\varphi_{\step}(S_1,S_2)=\mathrm{TRUE}$.  For the second conclusion we prove the
sharper statement that $\varphi_{\op}(S_1,S_2)=\mathrm{TRUE}$ for the operation
$\op$ whose disjunct of $\widehat{\varphi}_{\step,\bus}$ is satisfied; since
$\varphi_{\step}(S_1,S_2)$ is by definition the disjunction of the
$\varphi_{\op}(S_1,S_2)$ over all operations, the goal follows at once.
Every statement below is deterministic given
$\neg E_{\mathsf{pos}}\wedge\neg E_{\mathsf{upd}}$.
By Eq.~\eqref{eq:step-bus2}, the assumption expands as
\[
  \widehat{\varphi}_{\step}(\hat{S}_1,\hat{S}_2,\widehat{\BBB})
  \;=\;
  \widehat{\varphi}_{\step,\bus}(\hat{S}_1,\hat{S}_2,\widehat{\BBB})
  \;\wedge\;
  \widehat{\varphi}_{\keccak}(\widehat{\BBB})
  \;\wedge\;
  \widehat{\varphi}_{\poseidon}(\widehat{\BBB})
  \;\wedge\;
  \widehat{\varphi}_{\range}(\widehat{\BBB})
  \;=\;\mathrm{TRUE},
\]
so in particular
$\widehat{\varphi}_{\step,\bus}(\hat{S}_1,\hat{S}_2,\widehat{\BBB})=\mathrm{TRUE}$
and all chip predicates hold.

\smallskip
\noindent\emph{(No full-memory bus.)} The argument goes directly from the
committed, bus-parameterised predicate to the bus-free, full-memory predicate
$\varphi_{\step}(S_1,S_2)$; it never constructs a full-memory bus $\mathcal{B}$
and never invokes the bus-parameterised
form~\eqref{eq:step-expanded} of the step predicate.  This is deliberate, and
necessary: $\widehat{\BBB}$ is the bus of a whole segment
(Lemma~\ref{lem:segment}), so it carries range-check and precompile entries
belonging to \emph{other} steps, for which no memory vectors are extracted and
no full-memory states exist.  A predicate such as
$\varphi_{\keccak}(\mathcal{B})$, which quantifies over all entries of the bus,
therefore could not even be stated here, let alone established.  Instead each
bus-deferred obligation enters the case analysis on the committed side only, via
the chip predicates $\widehat{\varphi}_{\range}(\widehat{\BBB})$ and
$\widehat{\varphi}_{\op}(\widehat{\BBB})$ above, and is consumed at the single
bus entry belonging to this step.

\smallskip
\noindent\emph{(No commitment-injectivity argument.)} Because $\mem_2$ is
fixed by the reconstruction rule~\eqref{eq:mem-reconstruct} rather than chosen
by $\mathcal{A}$, the case analysis never has to infer an equality of vectors
from an equality of their commitments: each case obtains the memory component
of $\varphi_{\op}(S_1,S_2)$ directly from~\eqref{eq:mem-reconstruct}, and
obtains the invariant $\hat{S}_2.\mem=\mathsf{Commit}(\mem_2)$ either from the
commitment-equality conjunct of the disjunct (all cases but $\writ$) or from
update binding (the case $\writ$).  Position binding is used only to pin the
\emph{value} at the accessed address, at the single position $\addr$.
\begin{instruction}
  Here, and in some other places of the proofs, we implicitly used correctness of the vector commitment scheme, which we never explicitly defined. Team are expected to be less sloppy in this regard.
\end{instruction}

\smallskip

We case-split on the satisfied disjunct of $\widehat{\varphi}_{\step,\bus}$
(cf.~\eqref{eq:step-bus2}).  By Remark~\ref{rem:mem-ext-hyp4} the satisfied
disjunct is the one belonging to $\op=\code[\hat{S}_1.\pc]$, so the case
distinction below is the same one that drives the reconstruction
rule~\eqref{eq:mem-reconstruct}.

\smallskip
\noindent\emph{(Invariant, all cases but $\writ$.)} In the four non-write cases
$\op\neq\writ$, so~\eqref{eq:mem-reconstruct} sets $\mem_2:=\mem_1$, and the
satisfied disjunct asserts the commitment equality
$\hat{S}_1.\mem=\hat{S}_2.\mem$.  The invariant is then immediate,
\[
  \hat{S}_2.\mem=\hat{S}_1.\mem=\mathsf{Commit}(\mem_1)=\mathsf{Commit}(\mem_2),
\]
with no appeal to either binding property, and the memory conjunct
$\mem_2=\mem_1$ of $\varphi_{\op}$ holds by construction.  We record this once
and do not repeat it in the four cases below; only the write case, where the
commitment genuinely changes, needs update binding.

\smallskip
\begin{itemize}
\item \textbf{Non-memory arithmetic} ($\op_1,\ldots,\op_{10}$).
  The disjunct asserts
  $\widehat{\varphi}_{\op}(\hat{S}_1,\hat{S}_2)\wedge\hat{S}_1.\mem=\hat{S}_2.\mem$.
  The arithmetic part gives
  $\varphi'_{\op}(\hat{S}_1.\pc,\hat{S}_1.\regs,\hat{S}_2.\pc,\hat{S}_2.\regs)=\mathrm{TRUE}$,
  and $\mem_2=\mem_1$ by~\eqref{eq:mem-reconstruct}. Therefore
  $\varphi_{\op}(S_1,S_2)=\varphi'_{\op}(S_1,S_2)\wedge(\mem_2=\mem_1)=\mathrm{TRUE}$.

\item \textbf{Arithmetic with range checks} ($\op_{11},\ldots,\op_{20}$).
  The arithmetic part
\[
  \widehat{\varphi}_{\op}'(\hat{S}_1,\hat{S}_2)
\]
is
  memory-free, so $\varphi_{\op}'(S_1,S_2)=\mathrm{TRUE}$. The range-check
  requirement $(\range,\hat{S}_1)\in\widehat{\BBB}$ and
\[
  \widehat{\varphi}_{\range}(\widehat{\BBB})=\mathrm{TRUE}
\]
together give
  $\varphi_{\range}(S_1)=\mathrm{TRUE}$ (range checks depend only on registers,
  not memory), and $\mem_2=\mem_1$ by~\eqref{eq:mem-reconstruct}. Hence
\[
  \varphi_{\op}(S_1,S_2)=\varphi_{\op}'(S_1,S_2)\wedge\varphi_{\range}(S_1)\wedge(\mem_2=\mem_1)=\mathrm{TRUE}.
\]


\item \textbf{Read} ($\rea$).
  Here $\pi^{\mathsf{mem}}=\pi$ with $\hat{S}_1.\mem=\hat{S}_2.\mem$ and
\[
  \mathsf{Verify}(\hat{S}_1.\mem,\hat{S}_1.\regs[0],\hat{S}_2.\regs[1],\pi)=1.
\]
  Write $\addr:=\hat{S}_1.\regs[0]$ and $x:=\hat{S}_2.\regs[1]$. The honest
  opening $\mathsf{Open}(\mem_1,\addr)$ verifies against $\hat{S}_1.\mem$ to
  $\mem_1[\addr]$; were $\mem_1[\addr]\neq x$, this and $\pi$ would witness
  $E_{\mathsf{pos}}$ at $(R,p)=(\hat{S}_1.\mem,\addr)$, which is excluded. So
  $\mem_1[\addr]=x$, which is the memory-access conjunct
  $\mem_1[\regs_1[0]]=\regs_2[1]$ of~\eqref{eq:phi-read-decomp}.  The
  memory-preservation conjunct $\mem_2=\mem_1$
  of~\eqref{eq:phi-read-decomp} holds by~\eqref{eq:mem-reconstruct}.  The
  arithmetic part satisfies
  \[
    \varphi'_{\rea}(\hat{S}_1.\pc,\hat{S}_1.\regs,\hat{S}_2.\pc,\hat{S}_2.\regs)=\mathrm{TRUE}
  \]
  by Eq.~\eqref{eq:op-mem-comm-read}.  All three conjuncts hold, so
  $\varphi_{\rea}(S_1,S_2)=\mathrm{TRUE}$ per Eq.~\eqref{eq:phi-read-decomp}.

\item \textbf{Write} ($\writ$).
  Here $\pi^{\mathsf{mem}}=(\pi_{\writ},v_{\mathsf{old}})$ with
  \begin{align*}
    &\mathsf{Verify}(\hat{S}_1.\mem,\hat{S}_1.\regs[0],v_{\mathsf{old}},\pi_{\writ})=1,\\
    &\mathsf{Verify}(\hat{S}_2.\mem,\hat{S}_1.\regs[0],\hat{S}_1.\regs[1],\pi_{\writ})=1.
  \end{align*}
  Write $\addr:=\hat{S}_1.\regs[0]$ and $x:=\hat{S}_1.\regs[1]$, so that
  $\mem_2=\mem_1[\addr\mapsto x]$ by~\eqref{eq:mem-reconstruct}. The honest
  opening of $\mem_1$ at $\addr$ verifies against $\hat{S}_1.\mem$ to
  $\mem_1[\addr]$; were $\mem_1[\addr]\neq v_{\mathsf{old}}$, this and
  $\pi_{\writ}$ would witness $E_{\mathsf{pos}}$ under $\hat{S}_1.\mem$. So off
  $E_{\mathsf{pos}}$, $\mem_1[\addr]=v_{\mathsf{old}}$; that is, the shared path
  $\pi_{\writ}$ opens the honest commitment $\hat{S}_1.\mem=\mathsf{Commit}(\mem_1)$
  at $\addr$ to $\mem_1[\addr]$ and opens $\hat{S}_2.\mem$ at $\addr$ to $x$. Off
  $E_{\mathsf{upd}}$ this forces
  \[
    \hat{S}_2.\mem=\mathsf{Commit}(\mem_1[\addr\mapsto x])=\mathsf{Commit}(\mem_2),
  \]
  which is exactly the invariant for this case --- the one place in the proof
  where update binding is used, and the reason $\Com_{\mathsf{mem}}$ is required
  to have that property at all.  The memory conjunct
  $\mem_2=\mem_1[\regs_1[0]\mapsto\regs_1[1]]$
  of~\eqref{eq:phi-write-decomp} holds by~\eqref{eq:mem-reconstruct}, and with
  the arithmetic part
  \[
    \varphi'_{\writ}(\hat{S}_1.\pc,\hat{S}_1.\regs,\hat{S}_2.\pc,\hat{S}_2.\regs)=\mathrm{TRUE}
  \]
  from Eq.~\eqref{eq:op-mem-comm-write} we get
  $\varphi_{\writ}(S_1,S_2)=\mathrm{TRUE}$ per Eq.~\eqref{eq:phi-write-decomp}.

\item \textbf{Keccak/Poseidon precompile.}
  The disjunct of $\widehat{\varphi}_{\step,\bus}$ asserts
  $(\op,\hat{S}_1,\hat{S}_2)\in\widehat{\BBB}$. Since
  $\widehat{\varphi}_{\op}(\widehat{\BBB})=\mathrm{TRUE}$, the entry
  $(\op,\hat{S}_1,\hat{S}_2)$ satisfies
  $\widehat{\varphi}_{\op}(\hat{S}_1,\hat{S}_2)=\mathrm{TRUE}$, where
  $\varphi_{\op}$ (for $\op\in\{\keccak,\poseidon\}$) is the predicate defined in
  Section~\ref{sec:proof-arch} (``Predicates for individual operations'')
  asserting correctness of the precompile call.  Like every non-memory
  operation, $\varphi_{\op}$ decomposes as
  $\varphi'_{\op}(\pc_1,\regs_1,\pc_2,\regs_2)\wedge(\mem_2=\mem_1)$, and its
  pc-and-register part $\varphi'_{\op}$ --- the part that actually asserts
  correctness of the precompile call --- is memory-free and hence literally
  identical to $\widehat{\varphi}'_{\op}$.  We therefore obtain
  $\varphi'_{\op}(S_1,S_2)=\mathrm{TRUE}$.  The remaining conjunct is
  $\mem_2=\mem_1$, which holds by~\eqref{eq:mem-reconstruct}, as expected for a
  precompile that does not modify main memory.  Hence
  $\varphi_{\op}(S_1,S_2)=\varphi'_{\op}(S_1,S_2)\wedge(\mem_2=\mem_1)=\mathrm{TRUE}$.
\end{itemize}

In every case the invariant $\hat{S}_2.\mem=\mathsf{Commit}(\mem_2)$ holds, and
the bus-free, full-memory predicate
$\varphi_{\op}(S_1,S_2)=\mathrm{TRUE}$ holds for the operation $\op$ selected by
the satisfied disjunct.  Since $\varphi_{\step}(S_1,S_2)$ is the disjunction of
the $\varphi_{\op}(S_1,S_2)$ over all operations, the latter gives
$\varphi_{\step}(S_1,S_2)=\mathrm{TRUE}$, and both goals of Step~A are
established.  Note that the bus and the commitments have now disappeared from
the second conclusion, so no separate bus-elimination step is required.

\paragraph{Conclusion.}
Suppose $\mathcal{A}$ wins, i.e.\ all four hypotheses hold; in particular
$\widehat{\varphi}_{\step}(\hat{S}_1,\hat{S}_2,\widehat{\BBB})=\mathrm{TRUE}$ and
at least one of $\hat{S}_2.\mem=\mathsf{Commit}(\mem_2)$,
$\varphi_{\step}(S_1,S_2)=\mathrm{TRUE}$ fails. If neither $E_{\mathsf{pos}}$ nor
$E_{\mathsf{upd}}$ occurred, Step~A would give both, a contradiction; hence at
least one bad event occurs whenever $\mathcal{A}$ wins. By a union bound,
\[
  \varepsilon \;\le\; \Pr[E_{\mathsf{pos}}] + \Pr[E_{\mathsf{upd}}]
  \;\le\;
  \mathsf{Adv}^{\mathsf{pos}}_{\Com_{\mathsf{mem}}}(\mathcal{D}^{\mathsf{pos}})
  + \mathsf{Adv}^{\mathsf{upd}}_{\Com_{\mathsf{mem}}}(\mathcal{D}^{\mathsf{upd}}).
\]
\end{proof}

\begin{remark}[Memory reconstruction invariant]%
  \label{rem:mem-inheritance}
In Theorem~\ref{thm:main} Step~6 the memory vectors $\mem_k$ are reconstructed
inductively: $\mem_0$ is read off the initial state, and $\mem_{k+1}$ is
obtained from $\mem_k$ by the reconstruction rule~\eqref{eq:mem-reconstruct}
applied to the step $(\hat{S}_k,\hat{S}_{k+1})$.  The claim carried along the
induction is the \emph{commitment invariant}
\[
  \hat{S}_k.\mem=\mathsf{Commit}(\mem_k)\qquad\text{for all }k.
\]

The invariant is \emph{not} an independent observation requiring its own
appeal to the two binding properties: it is the first conclusion of
Proposition~\ref{prop:memory-extractability}.  The base case $k=0$ holds by
construction, and the inductive step is one invocation of the proposition,
which --- from $\hat{S}_k.\mem=\mathsf{Commit}(\mem_k)$ and
$\widehat{\varphi}_{\step}(\hat{S}_k,\hat{S}_{k+1},\widehat{\BBB})=\mathrm{TRUE}$
--- returns both $\hat{S}_{k+1}.\mem=\mathsf{Commit}(\mem_{k+1})$, which
carries the induction, and $\varphi_{\step}(S_k,S_{k+1})=\mathrm{TRUE}$, which
is the per-step obligation of~\eqref{eq:relation-star}.  Both cost the same
single pair of binding events, so the two guarantees are obtained together
rather than paid for twice.  Within the proposition the work splits as follows:
\begin{itemize}
  \item \textbf{Write step} at address $\addr$ with new value $x$: the
    commitment genuinely changes, and only here is update binding used.
    Position binding first identifies $v_{\mathsf{old}}$ with $\mem_k[\addr]$,
    which turns $\pi_{\writ}$ into an opening of the \emph{honest} commitment
    $\hat{S}_k.\mem=\mathsf{Commit}(\mem_k)$; update binding then forces
    $\hat{S}_{k+1}.\mem=\mathsf{Commit}(\mem_k[\addr\mapsto x])
    =\mathsf{Commit}(\mem_{k+1})$.
  \item \textbf{All other steps}: $\mem_{k+1}=\mem_k$
    by~\eqref{eq:mem-reconstruct} and $\hat{S}_{k+1}.\mem=\hat{S}_k.\mem$
    (enforced by the memory-equality conjunct of $\widehat{\varphi}_{op}$
    in~\eqref{eq:step-bus2}), so the invariant is inherited with no
    cryptographic assumption at all.
\end{itemize}

Position binding is additionally what pins the value returned by a read: for a
read step at index $k$, Eq.~\eqref{eq:op-mem-comm-read} provides $\pi$ with
\[
  \mathsf{Verify}(\hat{S}_k.\mem,\hat{S}_k.\regs[0],\hat{S}_{k+1}.\regs[1],\pi)=1,
\]
and together with the invariant $\hat{S}_k.\mem=\mathsf{Commit}(\mem_k)$ this
gives $\hat{S}_{k+1}.\regs[1]=\mem_k[\hat{S}_k.\regs[0]]$ --- the memory-access
conjunct of~\eqref{eq:phi-read-decomp}.  This is the read case of Step~A of the
proposition, stated here in the notation of Theorem~\ref{thm:main}.
\end{remark}

%----------------------------------------------------------------------
\subsection{Lemmas: descending the recursion tree}\label{sec:security:lemmas}
%----------------------------------------------------------------------

We state one lemma per layer of the recursion tree, proceeding from the
outermost (\texttt{embed}) circuit inward to the inner-circuit layer.  Each
lemma composes the straight-line extractors guaranteed by
Definition~\ref{def:extractable}.

\paragraph{Advantage notation.}
We collect the advantage notation introduced with the definitions above:
$\mathsf{Adv}^{\mathsf{ks}}_{\Pi}$ (knowledge-soundness error,
Definition~\ref{def:extractable});
$\mathsf{Adv}^{\mathsf{pos}}_{\Com}$ and $\mathsf{Adv}^{\mathsf{upd}}_{\Com}$
(position-binding and update-binding advantages for the memory commitment
scheme $\Com_{\mathsf{mem}}$, Definition~\ref{def:binding}); and $
  \mathsf{Adv}^{\mathsf{cr}}_{\Com_{\mathsf{bus}}}$
(collision-resistance advantage
for the bus hash function, Definition~\ref{def:bus-cr}).
All advantages are functions of the security parameter $\lambda$.
In each lemma below we exhibit explicit reduction adversaries $\mathcal{D}$
and state their running times.
When a bound has the form $k\cdot\mathsf{Adv}(\mathcal{D})$, the reduction
$\mathcal{D}$ denotes the maximum-advantage adversary among the $k$ per-node
(or per-step) reductions constructed in the proof; each such reduction runs in
\[
  \mathsf{Time}(\mathcal{A})+\mathsf{poly}(\lambda).
\]


% ---- Lemma 1: embed -------------------------------------------------
\subsubsection{Embed layer: base of the composition}\label{sec:lem:embed}

\begin{definition}[Embed extraction game]\label{game:embed}
For a zkVM system, an extractor $\mathcal{E}_4$, and a PPT adversary
$\mathcal{A}$, define the game
$\mathsf{Exp}^{\mathsf{embed}}_{\mathcal{E}_4,\mathcal{A}}(\lambda)$:
\begin{enumerate}
  \item \textbf{Adversary.} $\mathcal{A}(1^\lambda)$ outputs a statement
    $(\hat{S}_0,\hat{S}_T)$ and a proof $\pi^4$; the step count $T$ is a
    fixed system parameter.
  \item \textbf{Filter.} If
    $\Pi_4.\mathsf{Verify}((\hat{S}_0,\hat{S}_T),\pi^4)=0$, return $0$.
  \item \textbf{Extraction.} Run
    $\pi^3 \leftarrow \mathcal{E}_4(\hat{S}_0,\hat{S}_T,\pi^4)$.
  \item \textbf{Winning condition.} Return $1$ iff
    $\Pi_3.\mathsf{Verify}((\hat{S}_0,\hat{S}_T,T),\pi^3)=0$, i.e.\ the extractor
    failed to produce a valid $\mathcal{R}_4$ witness despite an accepting
    \texttt{embed} proof.
\end{enumerate}
Write $\mathsf{Adv}^{\mathsf{embed}}_{\mathcal{E}_4}(\mathcal{A}) :=
\Pr[\mathsf{Exp}^{\mathsf{embed}}_{\mathcal{E}_4,\mathcal{A}}(\lambda)=1]$,
where the probability is over the random coins of $\mathcal{A}$ and of the
extractor $\mathcal{E}_4$ (the two $\mathsf{Verify}$ steps are deterministic);
as in Definition~\ref{def:extractable}, any setup/CRS or random-oracle
sampling is left implicit.
\end{definition}


\begin{lemma}[Embed extraction]\label{lem:embed}
Assume $\Pi_4$ is knowledge-sound (Definition~\ref{def:extractable}).
We construct an extractor $\mathcal{E}_4$ (given in the proof) such that for every PPT
adversary $\mathcal{A}$ there is a reduction adversary $\mathcal{D}_4$ with
\[
  \mathsf{Adv}^{\mathsf{embed}}_{\mathcal{E}_4}(\mathcal{A})
  \;\le\; \mathsf{Adv}^{\mathsf{ks}}_{\Pi_4}(\mathcal{D}_4),
\]
where
$\mathsf{Time}(\mathcal{D}_4)=\mathsf{Time}(\mathcal{A})+\mathsf{poly}(\lambda)$.
This is the base layer of the extractor composition; the bound is a direct
instantiation of the knowledge soundness of $\Pi_4$, stated as a game only for
uniformity with Lemmas~\ref{lem:combine}--\ref{lem:segment}.
\end{lemma}

\begin{proof}
Let $\mathcal{E}_{\Pi_4}$ be the knowledge-soundness extractor for $\Pi_4$
guaranteed by Definition~\ref{def:extractable}, and set
$\mathcal{E}_4 := \mathcal{E}_{\Pi_4}$ (the $\mathcal{R}_4$ witness is exactly
$\pi^3$). Define the reduction adversary $\mathcal{D}_4$, \emph{playing the
knowledge-soundness game of $\Pi_4$ for relation $\mathcal{R}_4$}, as follows:
\begin{enumerate}
  \item Run $((\hat{S}_0,\hat{S}_T),\pi^4)\leftarrow\mathcal{A}(1^\lambda)$.
  \item Output the statement $x:=(\hat{S}_0,\hat{S}_T)$ and proof $\pi^4$ to its
    own challenger.
\end{enumerate}
$\mathcal{D}_4$ makes a single invocation of $\mathcal{A}$ and no other
computation, so
$\mathsf{Time}(\mathcal{D}_4)=\mathsf{Time}(\mathcal{A})+\mathsf{poly}(\lambda)$.

The knowledge-soundness game runs $\mathcal{E}_{\Pi_4}(x,\pi^4)$ to obtain a
candidate witness, which by construction is the value $\pi^3$ returned by
$\mathcal{E}_4$. Comparing the two games: both accept the same $(x,\pi^4)$ at
the filter step, and $\mathsf{Exp}^{\mathsf{embed}}$ returns $1$ exactly when
$\pi^4$ verifies but $(x;\pi^3)\notin\mathcal{R}_4$ --- precisely the event that
$\mathcal{D}_4$ wins the knowledge-soundness game of $\Pi_4$. Hence
$\mathsf{Adv}^{\mathsf{embed}}_{\mathcal{E}_4}(\mathcal{A})
=\mathsf{Adv}^{\mathsf{ks}}_{\Pi_4}(\mathcal{D}_4)$, giving the stated bound.
\end{proof}

% ---- Lemma 2: combine -----------------------------------------------
\subsubsection{Combine layer: unrolling the recursion tree}\label{sec:lem:combine}

We establish well-foundedness of the recursion tree.

\begin{remark}[Well-foundedness of \texttt{combine}]\label{rem:wellfounded}
By convention, every \texttt{combine} statement $(\hat{S}_0^L, \hat{S}_N^R, N)$
carries its step count $N$ as a public input.
The relation $\mathcal{R}_3$ enforces $N_{{\seg}} \mid N_L$, $N_{{\seg}} \mid N_R$,
$N_L \geq N_{{\seg}}$, $N_R \geq N_{{\seg}}$, and $N_L+N_R=N$, so each child's
step count is strictly smaller than its parent's.
Hence the recursion tree is finite and strong induction on $N$ is well-founded.
\end{remark}

Throughout this lemma we assume $2N_{\seg} \leq N \leq T$ and $T \leq \mathsf{poly}(\lambda)$
(the step count is a system parameter, Definition~\ref{def:relation-star}), so $m = N/N_{\seg}$ is polynomially bounded.

\begin{definition}[Combine extraction game]\label{game:combine}
For an extractor $\mathcal{E}_3$, a step count $N$ with $N_{\seg}\mid N$ (write
$m:=N/N_{\seg}$), and PPT $\mathcal{A}$, define the game
$\mathsf{Exp}^{\mathsf{comb}}_{N,\mathcal{E}_3,\mathcal{A}}(\lambda)$:
\begin{enumerate}
  \item \textbf{Adversary.} $\mathcal{A}(1^\lambda)$ outputs
    $(\hat{S}_0,\hat{S}_N,N)$ and a proof $\pi^3$.
  \item \textbf{Filter.} If $\Pi_3.\mathsf{Verify}((\hat{S}_0,\hat{S}_N,N),\pi^3)=0$,
    return $0$.
  \item \textbf{Extraction.} Run
\[
  (\hat{S}^{(0)},\ldots,\hat{S}^{(m)};\,\pi^2_1,\ldots,\pi^2_m)
    \leftarrow\mathcal{E}_3(\hat{S}_0,\hat{S}_N,N,\pi^3).
\]
  \item \textbf{Winning condition.} Return $1$ unless
    $\hat{S}^{(0)}=\hat{S}_0$, $\hat{S}^{(m)}=\hat{S}_N$, and
\[
  \Pi_2.\mathsf{Verify}((\hat{S}^{(i-1)},\hat{S}^{(i)},N_{\seg}),\pi^2_i)=1
\]
for all $i\in[m]$.
\end{enumerate}
Write
\[
  \mathsf{Adv}^{\mathsf{comb}}_{N,\mathcal{E}_3}(\mathcal{A}) :=
\Pr[\mathsf{Exp}^{\mathsf{comb}}_{N,\mathcal{E}_3,\mathcal{A}}(\lambda)=1],
\]
the probability taken as in Definition~\ref{game:embed}.
\end{definition}

\begin{lemma}[Combine tree unrolling]\label{lem:combine}
Assume $\Pi_3$ is knowledge-sound (Definition~\ref{def:extractable}), and let
$N\ge 2N_{\seg}$ with $N_{\seg}\mid N$ and $m:=N/N_{\seg}$. We construct an
extractor $\mathcal{E}_3$ (given in the proof) such that for every PPT $\mathcal{A}$ there is a
reduction adversary $\mathcal{D}_3$, running in
$\mathsf{Time}(\mathcal{A})+\mathsf{poly}(\lambda)$, with
\[
  \mathsf{Adv}^{\mathsf{comb}}_{N,\mathcal{E}_3}(\mathcal{A})
  \le (m-1)\cdot\mathsf{Adv}^{\mathsf{ks}}_{\Pi_3}(\mathcal{D}_3).
\]
\end{lemma}

\begin{proof}
The combine game returns $0$ unless $\mathcal{A}$'s \texttt{combine} proof
verifies, so we assume it does and bound the probability that extraction
nonetheless violates the winning condition; write $m:=N/N_{\seg}$. The argument
runs in four steps.

\begin{enumerate}[label=\textbf{Step~\arabic*.},leftmargin=*,itemsep=5pt]
\item \emph{Recursion tree and bad events.}
  The extractor $\mathcal{E}_3$ unrolls a binary tree in a fixed depth-first
  order: at each \texttt{combine} node it runs $\mathcal{E}_{\Pi_3}$ on that
  node's statement and proof to obtain a candidate witness
  $(\pi^L,\pi^R,\hat{S}^{mid},N_L,N_R)$, recurses into any child whose step
  count is at least $2N_{\seg}$ (a \texttt{combine} child, by the typing
  of~\eqref{eq:rel-combine-typed}), and stops at \texttt{convert} children
  ($N_{child}=N_{\seg}$).  For $t\ge 1$ define
  \[
    B_t:\quad\parbox[t]{0.8\linewidth}{the \emph{first} $\mathcal{E}_{\Pi_3}$
    invocation (in unroll order) that does not return a valid
    $\mathcal{R}_3$-witness for its node's statement is invocation number
    $t$.}
  \]
  The events $B_t$ are disjoint, and at every $B_t$ the invocation's proof
  verifies: the root's by the game's filter, a child's by its parent's valid
  witness.  At the moment of the first failure every earlier invocation
  returned a valid witness; the step counts of a valid witness partition the
  node's count into multiples of $N_{\seg}$, so the valid prefix contains at
  most $m-1$ internal nodes.  Moreover, if all $m-1$ possible internal
  invocations are valid, every leaf-child has step count exactly $N_{\seg}$
  and is a \texttt{convert} child by the typing
  of~\eqref{eq:rel-combine-typed}, so no further $\mathcal{E}_{\Pi_3}$
  invocation exists.  Hence only $B_1,\ldots,B_{m-1}$ are possible.  Off all
  of them (i.e.\ if no invocation fails), the tree has exactly $m-1$ internal
  (\texttt{combine}) nodes and $m$ leaves (\texttt{convert} proofs).

\item \emph{One reduction per invocation.}
  $\mathcal{D}_3^{(t)}$ runs $\mathcal{A}(1^\lambda)$ and unrolls the tree up
  to the $t$-th $\mathcal{E}_{\Pi_3}$ invocation (deterministic given
  $\mathcal{A}$'s output and the extractions at the earlier invocations), then
  forwards that invocation's statement and proof to the $\Pi_3$
  knowledge-soundness challenger. It wins whenever $B_t$ occurs, so
  \[
    \Pr[B_t]\le\mathsf{Adv}^{\mathsf{ks}}_{\Pi_3}(\mathcal{D}_3^{(t)}).
  \]
  Each $\mathcal{D}_3^{(t)}$ makes a single run of $\mathcal{A}$ and runs in
  $\mathsf{Time}(\mathcal{A})+\mathsf{poly}(\lambda)$.

\item \emph{Off the bad events, the extractor succeeds.}
  Assume $\neg B_1\wedge\cdots\wedge\neg B_{m-1}$; every statement in this step
  is then deterministic. We argue by strong induction on $N$ (justified by
  Remark~\ref{rem:wellfounded}).
  \begin{enumerate}[label=(\alph*),leftmargin=*,topsep=3pt]
  \item \emph{Base case $N=2N_{\seg}$ ($m=2$, one internal node).}
    Off $B_1$ (the root is the only $\mathcal{E}_{\Pi_3}$ invocation),
    applying $\mathcal{E}_{\Pi_3}$ to the root yields a valid witness with
    $N_L+N_R=2N_{\seg}$ and $N_{\seg}\mid N_L$, $N_{\seg}\mid N_R$, forcing
    $N_L=N_R=N_{\seg}$.  By the step-count typing
    of~\eqref{eq:rel-combine-typed} both children are verified under
    $\Vfy_2$, i.e.\ they are the two leaf \texttt{convert} proofs; the
    extractor outputs them together with the boundary states
    $(\hat{S}_0,\hat{S}^{mid},\hat{S}_N)$, meeting the winning condition.
    (No $\mathcal{R}_3$-witness exists at step count $N_{\seg}$ --- the
    typing forces $N_L,N_R\ge N_{\seg}$, hence $N_L+N_R\ge 2N_{\seg}$ ---
    which is why this lemma, like the system parameters of
    Definition~\ref{def:relation-star}, requires $N\ge 2N_{\seg}$: a \texttt{combine}
    statement at step count $N_{\seg}$ is unsatisfiable, yet a proof for it
    could still \emph{verify}, and no extractor could then meet the winning
    condition.)

  \item \emph{Inductive step $N>2N_{\seg}$.}
    Let $\mathcal{A}$ output $\pi^3$ for $(\hat{S}_0,\hat{S}_N,N)$. Off
    $B_1$ (the root invocation), applying $\mathcal{E}_{\Pi_3}$ to
    $(\hat{S}_0,\hat{S}_N,N,\pi^3)$ yields a valid witness
    \[
      (\pi^L,\pi^R,\hat{S}^{mid},N_L,N_R)
    \]
    for the statement $(\hat{S}_0,\hat{S}_N,N)$, i.e.\ with $N_L+N_R=N$,
    $N_{\seg}\mid N_L$, $N_{\seg}\mid N_R$, and each child verified according
    to its step count~\eqref{eq:rel-combine-typed}:
    \begin{align*}
      &\bigl(N_L=N_{\seg}\wedge\Vfy_2((\hat{S}_0,\hat{S}^{mid},N_L),\pi^L)=1\bigr)\\
      &\quad\lor\;
       \bigl(N_L\ge 2N_{\seg}\wedge\Vfy_3((\hat{S}_0,\hat{S}^{mid},N_L),\pi^L)=1\bigr), \\
      &\bigl(N_R=N_{\seg}\wedge\Vfy_2((\hat{S}^{mid},\hat{S}_N,N_R),\pi^R)=1\bigr)\\
      &\quad\lor\;
       \bigl(N_R\ge 2N_{\seg}\wedge\Vfy_3((\hat{S}^{mid},\hat{S}_N,N_R),\pi^R)=1\bigr).
    \end{align*}
    \emph{Left child.} If $N_L=N_{\seg}$, the typing
    of~\eqref{eq:rel-combine-typed} gives $\Vfy_2(\cdot,\pi^L)=1$, and $\pi^L$
    is a leaf \texttt{convert} proof $\pi^2_L:=\pi^L$ covering
    $[\hat{S}_0,\hat{S}^{mid}]$. If instead $N_L\ge 2N_{\seg}$, the typing
    gives $\Vfy_3(\cdot,\pi^L)=1$, so $2N_{\seg}\le N_L<N$ and the left
    subtree is a smaller \texttt{combine} instance; its $\mathcal{E}_{\Pi_3}$
    invocations are among those indexed by the $\{B_t\}$,
    so off those events the induction hypothesis (applied to the extractor
    $\mathcal{E}_3^{(N_L)}$ on $(\hat{S}_0,\hat{S}^{mid},N_L,\pi^L)$) yields
    $m_L:=N_L/N_{\seg}$ \texttt{convert} proofs $\pi^2_1,\ldots,\pi^2_{m_L}$ and
    boundary states
    $\hat{S}^{(0)}=\hat{S}_0,\hat{S}^{(1)},\ldots,\hat{S}^{(m_L)}=\hat{S}^{mid}$.

    \emph{Right child.} Symmetrically, the range $[\hat{S}^{mid},\hat{S}_N]$ at
    step count $N_R<N$ yields $m_R:=N_R/N_{\seg}$ \texttt{convert} proofs.

    \emph{Concatenate.} Joining the two lists gives $m=m_L+m_R=N/N_{\seg}$
    \texttt{convert} proofs and a chain
    $\hat{S}^{(0)}=\hat{S}_0,\ldots,\hat{S}^{(m)}=\hat{S}_N$ with
    $\Pi_2.\mathsf{Verify}((\hat{S}^{(i-1)},\hat{S}^{(i)},N_{\seg}),\pi^2_i)=1$
    for all $i\in[m]$ --- exactly the winning condition of
    $\mathsf{Exp}^{\mathsf{comb}}$.
  \end{enumerate}

\item \emph{Conclusion (union bound).}
  Off all bad events the extractor meets the winning condition, so the game
  outputs $1$ only if some $B_t$ occurs, and only $B_1,\ldots,B_{m-1}$ are
  possible (Step~1), so by a union bound
  \[
    \mathsf{Adv}^{\mathsf{comb}}_{N,\mathcal{E}_3}(\mathcal{A})
    \;\le\;\sum_{t=1}^{m-1}\Pr[B_t]
    \;\le\;(m-1)\cdot\mathsf{Adv}^{\mathsf{ks}}_{\Pi_3}(\mathcal{D}_3),
  \]
  where $\mathcal{D}_3$ denotes the maximum-advantage reduction among
  $\mathcal{D}_3^{(1)},\ldots,\mathcal{D}_3^{(m-1)}$ (the convention recorded in
  the ``Advantage notation'' paragraph). Since $m=N/N_{\seg}$ is polynomially
  bounded (Definition~\ref{def:relation-star}), the right-hand side is negligible whenever
  $\mathsf{Adv}^{\mathsf{ks}}_{\Pi_3}$ is. The extractor invokes
  $\mathcal{E}_{\Pi_3}$ once per internal node, each call costing at most
  $\mathsf{Time}(\mathcal{A})+\mathsf{poly}(\lambda)$, so it runs in
  $(m-1)\cdot\mathsf{Time}(\mathcal{A})+\mathsf{poly}(\lambda)$.
\end{enumerate}
\end{proof}

% ---- Lemma 3: convert -----------------------------------------------
\subsubsection{Convert layer: normalization}\label{sec:lem:convert}

\begin{definition}[Convert extraction game]\label{game:convert}
For an extractor $\mathcal{E}_2$ and PPT $\mathcal{A}$, define the game
$\mathsf{Exp}^{\mathsf{conv}}_{\mathcal{E}_2,\mathcal{A}}(\lambda)$:
\begin{enumerate}
  \item \textbf{Adversary.} $\mathcal{A}(1^\lambda)$ outputs
    $(\hat{S}_{in},\hat{S}_{out},N_{\seg})$ and a proof $\pi^2$.
  \item \textbf{Filter.} If
    $\Pi_2.\mathsf{Verify}((\hat{S}_{in},\hat{S}_{out},N_{\seg}),\pi^2)=0$,
    return $0$.
  \item \textbf{Extraction.} Run
    $\pi^1\leftarrow\mathcal{E}_2(\hat{S}_{in},\hat{S}_{out},N_{\seg},\pi^2)$.
  \item \textbf{Winning condition.} Return $1$ iff
    $\Pi_1.\mathsf{Verify}((\hat{S}_{in},\hat{S}_{out}),\pi^1)=0$.
\end{enumerate}
Write $\mathsf{Adv}^{\mathsf{conv}}_{\mathcal{E}_2}(\mathcal{A}) :=
\Pr[\mathsf{Exp}^{\mathsf{conv}}_{\mathcal{E}_2,\mathcal{A}}(\lambda)=1]$,
the probability taken as in Definition~\ref{game:embed}.
\end{definition}

\begin{lemma}[Convert extraction]\label{lem:convert}
Assume $\Pi_2$ is knowledge-sound (Definition~\ref{def:extractable}).
We construct an extractor $\mathcal{E}_2$ (given in the proof) such that for every PPT
$\mathcal{A}$ there is a reduction adversary $\mathcal{D}_2$, running in
$\mathsf{Time}(\mathcal{A})+\mathsf{poly}(\lambda)$, with
\[
  \mathsf{Adv}^{\mathsf{conv}}_{\mathcal{E}_2}(\mathcal{A})
  \le \mathsf{Adv}^{\mathsf{ks}}_{\Pi_2}(\mathcal{D}_2).
\]
\end{lemma}

\begin{proof}
Set $\mathcal{E}_2:=\mathcal{E}_{\Pi_2}$ (the $\mathcal{R}_2$ witness is exactly
$\pi^1$), and let $\mathcal{D}_2$ run $\mathcal{A}$ and forward
\[
  ((\hat{S}_{in},\hat{S}_{out},N_{\seg}),\pi^2)
\]
to its challenger. The convert
game returns $1$ exactly when $\pi^2$ verifies but
\[
  \Pi_1.\mathsf{Verify}((\hat{S}_{in},\hat{S}_{out}),\pi^1)=0;
\]
this implies
$(x;\pi^1)\notin\mathcal{R}_2$, the event that $\mathcal{D}_2$ wins the
knowledge-soundness game of $\Pi_2$, so
$\mathsf{Adv}^{\mathsf{conv}}_{\mathcal{E}_2}(\mathcal{A})\le
\mathsf{Adv}^{\mathsf{ks}}_{\Pi_2}(\mathcal{D}_2)$.  (Equality can fail when
$\mathcal{A}$'s statement carries a step count other than $N_{\seg}$: the
statement is then outside the support of $\mathcal{R}_2$, so $\mathcal{D}_2$
wins whenever $\pi^2$ verifies, while the convert game may still return
$0$.)
\end{proof}

% ---- Lemma 4: segment -----------------------------------------------
\subsubsection{Segment layer: inner proofs to a segment}\label{sec:lem:segment}

\begin{definition}[Segment extraction game]\label{game:segment}
For an extractor $\mathcal{E}_1$ and PPT $\mathcal{A}$, define the game
$\mathsf{Exp}^{\mathsf{seg}}_{\mathcal{E}_1,\mathcal{A}}(\lambda)$:
\begin{enumerate}
  \item \textbf{Adversary.} $\mathcal{A}(1^\lambda)$ outputs
    $(\hat{S}_{in},\hat{S}_{out})$ and a proof $\pi^1$.
  \item \textbf{Filter.} If
    $\Pi_1.\mathsf{Verify}((\hat{S}_{in},\hat{S}_{out}),\pi^1)=0$, return $0$.
  \item \textbf{Extraction.} Run
    $(\widehat{\BBB},\hat{S}_1,\ldots,\hat{S}_{N_{\seg}+1},\{\pi^{\mathsf{mem}}_i\}_{i=1}^{N_{\seg}})
    \leftarrow\mathcal{E}_1(\hat{S}_{in},\hat{S}_{out},\pi^1)$.
  \item \textbf{Winning condition.} Return $1$ unless
    $\hat{S}_1=\hat{S}_{in}$, $\hat{S}_{N_{\seg}+1}=\hat{S}_{out}$, and
    \[
      \bigwedge_{i=1}^{N_{\seg}}\widehat{\varphi}_{\step}(\hat{S}_i,\hat{S}_{i+1},\widehat{\BBB},\pi^{\mathsf{mem}}_i)=\mathrm{TRUE}.
    \]
\end{enumerate}
Write $\mathsf{Adv}^{\mathsf{seg}}_{\mathcal{E}_1}(\mathcal{A}) :=
\Pr[\mathsf{Exp}^{\mathsf{seg}}_{\mathcal{E}_1,\mathcal{A}}(\lambda)=1]$,
the probability taken as in Definition~\ref{game:embed}.
\end{definition}

\begin{lemma}[Segment extraction]\label{lem:segment}
Assume $\Pi_1$ and each $\Pi_{0,j}$ ($j\in\{\step,\keccak,\poseidon,\range\}$)
are knowledge-sound (Definition~\ref{def:extractable}), and $\Com_{\mathsf{bus}}$
is collision-resistant (Definition~\ref{def:bus-cr}). We construct an extractor
$\mathcal{E}_1$ (given in the proof) such that for every PPT $\mathcal{A}$ there are reduction
adversaries
$\mathcal{D}_1,\mathcal{D}_{0,\step},\mathcal{D}_{0,\keccak},\mathcal{D}_{0,\poseidon},\mathcal{D}_{0,\range},\mathcal{D}_{\mathsf{bus}}$,
each running in $\mathsf{Time}(\mathcal{A})+\mathsf{poly}(\lambda)$, with
\begin{align*}
  \mathsf{Adv}^{\mathsf{seg}}_{\mathcal{E}_1}(\mathcal{A})
  \le{}& \mathsf{Adv}^{\mathsf{ks}}_{\Pi_1}(\mathcal{D}_1)
  + \mathsf{Adv}^{\mathsf{ks}}_{\Pi_{0,\step}}(\mathcal{D}_{0,\step})
  + \mathsf{Adv}^{\mathsf{ks}}_{\Pi_{0,\keccak}}(\mathcal{D}_{0,\keccak})\\
  &+ \mathsf{Adv}^{\mathsf{ks}}_{\Pi_{0,\poseidon}}(\mathcal{D}_{0,\poseidon})
  + \mathsf{Adv}^{\mathsf{ks}}_{\Pi_{0,\range}}(\mathcal{D}_{0,\range})
  + \mathsf{Adv}^{\mathsf{cr}}_{\Com_{\mathsf{bus}}}(\mathcal{D}_{\mathsf{bus}}).
\end{align*}
\end{lemma}

\begin{proof}
The segment game returns $0$ unless $\pi^1$ verifies, so we assume it does and
bound the probability that extraction violates the winning condition. We define
six bad events, one per reduction, and show that off all six the extractor meets
the winning condition; a union bound then gives the stated sum.

\paragraph{Bad events and their reductions.}
The six events below refer to the intermediate values produced by the
corresponding steps of the extractor $\mathcal{E}_1$, defined in the next
paragraph; on any fixed run, read the extractor's steps first and interpret
each event as ``that step's extraction fails''.
Each reduction runs $\mathcal{A}(1^\lambda)$ once (invoking the earlier partial
extractors as a subroutine) and runs in
$\mathsf{Time}(\mathcal{A})+\mathsf{poly}(\lambda)$.
\begin{itemize}
  \item $B_1$ (\emph{$\Pi_1$ extraction}): $\pi^1$ verifies but
    $\mathcal{E}_{\Pi_1}$ returns a witness
    $(C_{\hat{\mathcal{B}}},\pi_{\step},\pi_{\keccak},\pi_{\poseidon},\pi_{\range})\notin\mathcal{R}_1$.
    $\mathcal{D}_1$ forwards $((\hat{S}_{in},\hat{S}_{out}),\pi^1)$ to the $\Pi_1$
    challenger, so
    $\Pr[B_1]\le\mathsf{Adv}^{\mathsf{ks}}_{\Pi_1}(\mathcal{D}_1)$.
  \item $B_{0,\step}$ (\emph{\texttt{inner-step} extraction}): $\pi_{\step}$
    verifies but $\mathcal{E}_{0,\step}$ returns a witness
    $\notin\mathcal{R}_{0,\step}$. Here $\mathcal{D}_{0,\step}$ is the partial
    extractor running Step~1, forwarding $(x_{0,\step},\pi_{\step})$, so that
    $\Pr[B_{0,\step}]$ is at most
    $\mathsf{Adv}^{\mathsf{ks}}_{\Pi_{0,\step}}(\mathcal{D}_{0,\step})$.
  \item $B_{0,j}$ for $j\in\{\keccak,\poseidon,\range\}$ (\emph{inner-chip
    extraction}): $\pi_j$ verifies but $\mathcal{E}_{0,j}$ returns a witness
    $\notin\mathcal{R}_{0,j}$. $\mathcal{D}_{0,j}$ is the partial extractor
    running Steps~1--2, forwarding $(x_{0,j},\pi_j)$; so
    $\Pr[B_{0,j}]\le\mathsf{Adv}^{\mathsf{ks}}_{\Pi_{0,j}}(\mathcal{D}_{0,j})$.
  \item $B_{\mathsf{bus}}$ (\emph{bus collision}): the extracted buses
    $\widehat{\BBB}_{\step},\widehat{\BBB}_{\keccak},\widehat{\BBB}_{\poseidon},\widehat{\BBB}_{\range}$
    are not all equal, although all four commit to $C_{\hat{\mathcal{B}}}$.
    $\mathcal{D}_{\mathsf{bus}}$ outputs any unequal pair as a
    $\Com_{\mathsf{bus}}$ collision, so
    $\Pr[B_{\mathsf{bus}}]\le\mathsf{Adv}^{\mathsf{cr}}_{\Com_{\mathsf{bus}}}(\mathcal{D}_{\mathsf{bus}})$.
\end{itemize}

\paragraph{The extractor $\mathcal{E}_1$ (six steps).}

\textbf{Step~1.} Apply $\mathcal{E}_{\Pi_1}$ to $(\hat{S}_{in},\hat{S}_{out},\pi^1)$
to obtain
\[
  (C_{\hat{\mathcal{B}}},\pi_{\step},\pi_{\keccak},\pi_{\poseidon},\pi_{\range});
\]
off $B_1$ these satisfy
\[
  \Pi_{0,j}.\mathsf{Verify}(x_{0,j},\pi_j)=1
  \quad\text{for all }j\in\{\step,\keccak,\poseidon,\range\},
\]
with $x_{0,\step}=(\hat{S}_{in},\hat{S}_{out},C_{\hat{\mathcal{B}}})$ and
$x_{0,j}=C_{\hat{\mathcal{B}}}$ otherwise.

\textbf{Step~2.} Apply $\mathcal{E}_{0,\step}$ (Lemma~\ref{lem:inner}) to
$(x_{0,\step},\pi_{\step})$; off $B_{0,\step}$ obtain
$(\widehat{\BBB}_{\step},\{\hat{S}_i\},\{\pi^{\mathsf{mem}}_i\})$ with
$\hat{S}_1=\hat{S}_{in}$, $\hat{S}_{N_{\seg}+1}=\hat{S}_{out}$,
\[
  \textstyle\bigwedge_i\widehat{\varphi}_{\step,\bus}(\hat{S}_i,\hat{S}_{i+1},\widehat{\BBB}_{\step},\pi^{\mathsf{mem}}_i)=\mathrm{TRUE},
\]
and $\Com_{\mathsf{bus}}(\widehat{\BBB}_{\step})=C_{\hat{\mathcal{B}}}$.

\textbf{Step~3.} Apply $\mathcal{E}_{0,j}$ to $(x_{0,j},\pi_j)$ for
$j\in\{\keccak,\poseidon,\range\}$; off $B_{0,j}$ obtain buses
$\widehat{\BBB}_{\keccak},\widehat{\BBB}_{\poseidon},\widehat{\BBB}_{\range}$,
each satisfying its chip predicate and committing to $C_{\hat{\mathcal{B}}}$.

\textbf{Step~4.} Set $\widehat{\BBB}:=\widehat{\BBB}_{\step}$.

\textbf{Output.} $\mathcal{E}_1$ returns $\widehat{\BBB}$, the states
$\hat{S}_1,\ldots,\hat{S}_{N_{\seg}+1}$ of Step~2, and the memory-opening
witnesses $\{\pi^{\mathsf{mem}}_i\}_{i=1}^{N_{\seg}}$ (entries of the
$\relation_{0,\step}$ witness).

\paragraph{Off all six events, the extractor wins.}
Assume $\neg B_1\wedge\neg B_{0,\step}\wedge\bigwedge_{j}\neg B_{0,j}\wedge\neg B_{\mathsf{bus}}$.

\textbf{Bus consistency.} All four buses commit to $C_{\hat{\mathcal{B}}}$, and
off $B_{\mathsf{bus}}$ they are equal; this is consistent with the choice
$\widehat{\BBB}:=\widehat{\BBB}_{\step}$.

\textbf{Boundary conditions.} Step~2 gives $\hat{S}_1=\hat{S}_{in}$ and
$\hat{S}_{N_{\seg}+1}=\hat{S}_{out}$ directly from the $\mathcal{R}_{0,\step}$
witness.

\textbf{Step predicate.} Combining Step~2
($\widehat{\varphi}_{\step,\bus}(\hat{S}_i,\hat{S}_{i+1},\widehat{\BBB},\pi^{\mathsf{mem}}_i)=\mathrm{TRUE}$
for all $i$) with Step~3
($\widehat{\varphi}_{\keccak}(\widehat{\BBB})=\widehat{\varphi}_{\poseidon}(\widehat{\BBB})=\widehat{\varphi}_{\range}(\widehat{\BBB})=\mathrm{TRUE}$),
Eq.~\eqref{eq:step-bus2} gives
\begin{multline*}
  \widehat{\varphi}_{\step}(\hat{S}_i,\hat{S}_{i+1},\widehat{\BBB},\pi^{\mathsf{mem}}_i)
  \;=\;
  \widehat{\varphi}_{\step,\bus}(\hat{S}_i,\hat{S}_{i+1},\widehat{\BBB},\pi^{\mathsf{mem}}_i)\\
  \;\wedge\;
  \widehat{\varphi}_{\keccak}(\widehat{\BBB})
  \;\wedge\;
  \widehat{\varphi}_{\poseidon}(\widehat{\BBB})
  \;\wedge\;
  \widehat{\varphi}_{\range}(\widehat{\BBB})
  \;=\;\mathrm{TRUE}
\end{multline*}
for all $i\in\{1,\ldots,N_{\seg}\}$. Together with the boundary conditions this
is exactly the winning condition of $\mathsf{Exp}^{\mathsf{seg}}$.

\paragraph{Conclusion.}
Off all six bad events the game outputs $0$, so it outputs $1$ only if one of
them occurs. By a union bound,
\begin{align*}
  \mathsf{Adv}^{\mathsf{seg}}_{\mathcal{E}_1}(\mathcal{A})
  \le{}& \Pr[B_1]+\Pr[B_{0,\step}]+\sum_{j\in\{\keccak,\poseidon,\range\}}\Pr[B_{0,j}]+\Pr[B_{\mathsf{bus}}]\\
  \le{}& \mathsf{Adv}^{\mathsf{ks}}_{\Pi_1}(\mathcal{D}_1)
  + \mathsf{Adv}^{\mathsf{ks}}_{\Pi_{0,\step}}(\mathcal{D}_{0,\step})
  + \mathsf{Adv}^{\mathsf{ks}}_{\Pi_{0,\keccak}}(\mathcal{D}_{0,\keccak})\\
  &+ \mathsf{Adv}^{\mathsf{ks}}_{\Pi_{0,\poseidon}}(\mathcal{D}_{0,\poseidon})
  + \mathsf{Adv}^{\mathsf{ks}}_{\Pi_{0,\range}}(\mathcal{D}_{0,\range})
  + \mathsf{Adv}^{\mathsf{cr}}_{\Com_{\mathsf{bus}}}(\mathcal{D}_{\mathsf{bus}}),
\end{align*}
as claimed.
\end{proof}

% ---- Lemma 5: inner circuits ----------------------------------------
\subsubsection{Inner-circuit layer}\label{sec:lem:inner}

\begin{definition}[Inner-circuit extraction game]\label{game:inner}
For $j\in\{\step,\keccak,\poseidon,\range\}$, an extractor $\mathcal{E}_{0,j}$,
and PPT $\mathcal{A}$, define the game
$\mathsf{Exp}^{\mathsf{inner}}_{j,\mathcal{E}_{0,j},\mathcal{A}}(\lambda)$:
\begin{enumerate}
  \item \textbf{Adversary.} $\mathcal{A}(1^\lambda)$ outputs a public input
    $x_{0,j}$ and a proof $\pi_{0,j}$.
  \item \textbf{Filter.} If $\Pi_{0,j}.\mathsf{Verify}(x_{0,j},\pi_{0,j})=0$,
    return $0$.
  \item \textbf{Extraction.} Run
    $w_{0,j}\leftarrow\mathcal{E}_{0,j}(x_{0,j},\pi_{0,j})$.
  \item \textbf{Winning condition.} Return $1$ iff
    $(x_{0,j};w_{0,j})\notin\mathcal{R}_{0,j}$.
\end{enumerate}
Write $\mathsf{Adv}^{\mathsf{inner}}_{j,\mathcal{E}_{0,j}}(\mathcal{A}) :=
\Pr[\mathsf{Exp}^{\mathsf{inner}}_{j,\mathcal{E}_{0,j},\mathcal{A}}(\lambda)=1]$,
the probability taken as in Definition~\ref{game:embed}; this is precisely the
knowledge-soundness game of $\Pi_{0,j}$.
\end{definition}

\begin{lemma}[Inner-circuit extraction]\label{lem:inner}
Assume each $\Pi_{0,j}$, $j\in\{\step,\keccak,\poseidon,\range\}$, is
knowledge-sound (Definition~\ref{def:extractable}). For each $j$ we take the
extractor $\mathcal{E}_{0,j}:=\mathcal{E}_{\Pi_{0,j}}$ such that for
every PPT $\mathcal{A}$,
\[
  \mathsf{Adv}^{\mathsf{inner}}_{j,\mathcal{E}_{0,j}}(\mathcal{A})
  = \mathsf{Adv}^{\mathsf{ks}}_{\Pi_{0,j}}(\mathcal{D}_{0,j}),
  \qquad \mathcal{D}_{0,j}:=\mathcal{A},
\]
the game being that of knowledge soundness. The extracted witness $w_{0,j}$ has
the following structure:
\begin{itemize}
  \item $j=\step$:
    $w_{0,\step}=(\widehat{\BBB},\{\hat{S}_i\}_{1\le i\le N_{\seg}+1},
    \{\pi^{\mathsf{mem}}_i\}_{1\le i\le N_{\seg}})$
    with $\hat{S}_1=\hat{S}_{in}$,\; $\hat{S}_{N_{\seg}+1}=\hat{S}_{out}$,\;
\[
  \bigwedge_{i=1}^{N_{\seg}}\widehat{\varphi}_{\step,\bus}(\hat{S}_i,\hat{S}_{i+1},\widehat{\BBB},\pi^{\mathsf{mem}}_i)=\mathrm{TRUE},
\]
    and $\Com_{\mathsf{bus}}(\widehat{\BBB})=C_{\hat{\mathcal{B}}}$.
  \item $j\in\{\keccak,\poseidon,\range\}$: $w_{0,j}=\widehat{\BBB}$ with
    $\widehat{\varphi}_{j}(\widehat{\BBB})=\mathrm{TRUE}$ and
    $\Com_{\mathsf{bus}}(\widehat{\BBB})=C_{\hat{\mathcal{B}}}$.
\end{itemize}
\end{lemma}

\begin{proof}
$\mathcal{E}_{0,j}=\mathcal{E}_{\Pi_{0,j}}$ is the knowledge-soundness extractor
of $\Pi_{0,j}$, so $\mathsf{Exp}^{\mathsf{inner}}_{j}$ \emph{is} its
knowledge-soundness game and the advantages coincide; the extracted witness
lies in $\mathcal{R}_{0,j}$, giving the structure listed above.
\end{proof}

%----------------------------------------------------------------------
\subsection{Main theorem}\label{sec:security:thm}
%----------------------------------------------------------------------

\begin{theorem}[Correct-trace extractability]\label{thm:main}
Let $\mathsf{zkVM}$ be a zkVM system (Definition~\ref{def:zkvm}) with step count
$T$ a system parameter satisfying $N_{\seg}\mid T$, $T\ge 2N_{\seg}$, and
$T\le\mathsf{poly}(\lambda)$ (as in Definition~\ref{def:relation-star}),
and let $m := T/N_{\seg}$.
For every PPT adversary $\mathcal{A}$ with CTE advantage
$\varepsilon := \mathsf{Adv}^{\mathsf{cte}}_{\mathsf{zkVM}}(\mathcal{A})$
there exist PPT reduction adversaries
$\mathcal{D}_4$, $\mathcal{D}_3$, $\mathcal{D}_2$,
$\mathcal{D}_1$, $\{\mathcal{D}_{0,j}\}_{j}$, $\mathcal{D}_{\mathsf{bus}}$,
$\{\mathcal{D}^{\mathsf{pos}}_k\}_{k=1}^{T}$, $\{\mathcal{D}^{\mathsf{upd}}_k\}_{k=1}^{T}$,
$\mathcal{D}^{\mathsf{pos}}_{\mathsf{fin}}$
such that
\begin{align}\label{eq:cte-bound}
  \varepsilon &\leq
  \mathsf{Adv}^{\mathsf{ks}}_{\Pi_4}(\mathcal{D}_4)
  + (m{-}1)\,\mathsf{Adv}^{\mathsf{ks}}_{\Pi_3}(\mathcal{D}_3)
  + m\,\mathsf{Adv}^{\mathsf{ks}}_{\Pi_2}(\mathcal{D}_2) \nonumber\\
  &\quad
  + m\Bigl(
    \mathsf{Adv}^{\mathsf{ks}}_{\Pi_1}(\mathcal{D}_1)
    + \sum_{j\in\{\step,\keccak,\poseidon,\range\}}\mathsf{Adv}^{\mathsf{ks}}_{\Pi_{0,j}}(\mathcal{D}_{0,j})
    + \mathsf{Adv}^{\mathsf{cr}}_{\Com_{\mathsf{bus}}}(\mathcal{D}_{\mathsf{bus}})
  \Bigr) \nonumber\\
  &\quad
  + \sum_{k=1}^{T}\Bigl(
    \mathsf{Adv}^{\mathsf{pos}}_{\Com_{\mathsf{mem}}}(\mathcal{D}^{\mathsf{pos}}_k)
    + \mathsf{Adv}^{\mathsf{upd}}_{\Com_{\mathsf{mem}}}(\mathcal{D}^{\mathsf{upd}}_k)
  \Bigr)
  + \mathsf{Adv}^{\mathsf{pos}}_{\Com_{\mathsf{mem}}}(\mathcal{D}^{\mathsf{pos}}_{\mathsf{fin}}),
\end{align}
where the sum over $j$ runs over $\{\step,\keccak,\poseidon,\range\}$ and $k$ from $1$ to $T$;
the final summand accounts for the boundary check $\mem_T'=\mem_T$, a bad event
distinct from the $T$ per-step ones (Step~6 of the proof).
All reduction adversaries invoke $\mathcal{A}$ exactly once and run in time at
most $c\cdot\mathsf{Time}(\mathcal{A})+\mathsf{poly}(\lambda)$ with $c = 7m+1$
(computed in the running-time paragraph at the end of the proof).
In particular, if all advantage terms on the right-hand side are negligible in $\lambda$,
then $\mathsf{zkVM}$ is correct-trace extractable (Definition~\ref{def:cte}).
\end{theorem}

\begin{proof}
Let $\mathcal{A}$ be a PPT adversary with CTE advantage $\varepsilon$
(Definition~\ref{def:cte}); we bound $\varepsilon$.
We construct $\mathcal{E}_{\mathsf{zkVM}}$ by composing the sub-extractors
guaranteed by Lemmas~\ref{lem:embed}, \ref{lem:combine}, \ref{lem:convert},
\ref{lem:segment}, and~\ref{lem:inner}.  Each sub-extractor
is straight-line and is applied sequentially to the proof extracted by the
preceding layer.

\textbf{Step~1 (Embed layer).}
Apply $\mathcal{E}_4$ (Lemma~\ref{lem:embed}) to $(\hat{S}_0,\hat{S}_T,\pi^4)$
to obtain $\pi^3$ with
\[
  \Pi_3.\mathsf{Verify}((\hat{S}_0,\hat{S}_T,T),\pi^3)=1
\]
.
Fails with probability at most $\mathsf{Adv}^{\mathsf{ks}}_{\Pi_4}(\mathcal{D}_4)$,
where $\mathcal{D}_4$ runs in $\mathsf{Time}(\mathcal{A})+\mathsf{poly}(\lambda)$.%
\footnote{This composition is exactly where the over-idealization of
Remark~\ref{rem:idealized} and the ``Straight-line extraction'' remark is
invoked: for succinct $\Pi_4$ a straight-line extractor from $(x,\pi)$ alone
cannot exist, and composing straight-line extraction across the recursion would
require relativized SNARKs, which do not exist in the ROM. We idealize over
both, as declared there.}

\textbf{Step~2 (Combine layer).}
Apply $\mathcal{E}_3$ (Lemma~\ref{lem:combine}) to $(\hat{S}_0,\hat{S}_T,T,\pi^3)$
to obtain $m=T/N_{\seg}$ convert proofs $\pi^2_1,\ldots,\pi^2_m$ and boundary
states
\[
  \hat{S}^{(0)}=\hat{S}_0,\hat{S}^{(1)},\ldots,\hat{S}^{(m)}=\hat{S}_T
\]
.
Fails with probability at most
\[
  (m-1)\cdot\mathsf{Adv}^{\mathsf{ks}}_{\Pi_3}(\mathcal{D}_3)
\]
,
where $\mathcal{D}_3$ runs in
$\mathsf{Time}(\mathcal{A})+\mathsf{poly}(\lambda)$ and the extractor $\mathcal{E}_3$
in
\[
  (m-1)\cdot\mathsf{Time}(\mathcal{A})+\mathsf{poly}(\lambda)
\]
.

\textbf{Step~3 (Convert layer, applied $m$ times).}
For each $i\in[m]$, apply $\mathcal{E}_2$ (Lemma~\ref{lem:convert})
to $(\hat{S}^{(i-1)},\hat{S}^{(i)},N_{\seg},\pi^2_i)$.
Obtain $\pi^1_i$ with $\Pi_1.\mathsf{Verify}((\hat{S}^{(i-1)},\hat{S}^{(i)}),\pi^1_i)=1$.
Each invocation $i$ fails with probability at most
$\mathsf{Adv}^{\mathsf{ks}}_{\Pi_2}(\mathcal{D}_2)$; the total over $m$ calls is
$m\cdot\mathsf{Adv}^{\mathsf{ks}}_{\Pi_2}(\mathcal{D}_2)$.

\textbf{Step~4 (Segment layer, applied $m$ times).}
For each $i\in[m]$, apply $\mathcal{E}_1$ (Lemma~\ref{lem:segment},
which internally uses Lemma~\ref{lem:inner} and the collision resistance of $\Com_{\mathsf{bus}}$
(Definition~\ref{def:bus-cr})) to obtain
a bus $\widehat{\BBB}_i$ and states
$\hat{S}^{(i,0)},\ldots,\hat{S}^{(i,N_{\seg})}$ satisfying
\[
  \hat{S}^{(i,0)}=\hat{S}^{(i-1)},\quad
  \hat{S}^{(i,N_{\seg})}=\hat{S}^{(i)},\quad
  \bigwedge_{j=1}^{N_{\seg}}
    \widehat{\varphi}_{\step}(\hat{S}^{(i,j-1)},\hat{S}^{(i,j)},\widehat{\BBB}_i)
    =\mathrm{TRUE}.
\]
(The step predicates are evaluated on the extracted memory-opening witnesses
$\pi^{\mathsf{mem}}$, per the shorthand convention of
Section~\ref{sec:security:bus}.)
Each invocation $i$ fails with probability at most
\[
  \mathsf{Adv}^{\mathsf{ks}}_{\Pi_1}(\mathcal{D}_1)
+\sum_{j\in\{\step,\keccak,\poseidon,\range\}}\mathsf{Adv}^{\mathsf{ks}}_{\Pi_{0,j}}(\mathcal{D}_{0,j})
+\mathsf{Adv}^{\mathsf{cr}}_{\Com_{\mathsf{bus}}}(\mathcal{D}_{\mathsf{bus}});
\]
the total over $m$ calls is $m$ times that sum.

\textbf{Step~5 (Concatenation).}
Define $\hat{S}_k:=\hat{S}^{(\lfloor k/N_{\seg}\rfloor+1,\;k\bmod N_{\seg})}$
for $k=0,\ldots,T-1$, and $\hat{S}_T:=\hat{S}^{(m,N_{\seg})}$.  The boundary conditions $\hat{S}^{(i,N_{\seg})}=\hat{S}^{(i)}=\hat{S}^{(i+1,0)}$
ensure the chain is consistent across segment boundaries.  We obtain a
sequence $\hat{S}_0,\hat{S}_1,\ldots,\hat{S}_T$ of committed-memory states
with $\widehat{\varphi}_{\step}(\hat{S}_k,\hat{S}_{k+1},\widehat{\BBB}_{\lceil(k+1)/N_{\seg}\rceil})=\mathrm{TRUE}$
for all $k\in\{0,\ldots,T-1\}$.

\textbf{Step~6 (Memory reconstruction and memory extractability).}
Recall that CTE (Definition~\ref{def:cte}) requires the extractor to output a
witness for $\relation^*$~\eqref{eq:relation-star}, i.e.\ full-memory states
$S_0,\ldots,S_T$ with $\mem_0'=\mem_0$, $\mem_T'=\mem_T$,
and $\varphi_{\step}(S_k,S_{k+1})=\mathrm{TRUE}$ for all $k$.
Since the extractor receives the statement $(S_0,S_T)$, hence the full boundary
states (Remark~\ref{rem:cte-experiment}, Step~3), the initial memory $\mem_0$ is given
directly, so $\mem_0'=\mem_0$; the initial committed state satisfies
$\hat{S}_0.\mem=\mathsf{Commit}(\mem_0)$ by hypothesis.
For each step $k\in\{0,\ldots,T-1\}$, Lemma~\ref{lem:segment}
additionally yields the memory-opening witness $\pi^{\mathsf{mem}}_k$ for that
step; on a read it is parsed as the opening proof $\pi$ required by
$\widehat{\pred{}}_{\rea}$~\eqref{eq:op-mem-comm-read}, on a write as the pair
$(\pi_{\writ},v_{\mathsf{old}})$ required by
$\widehat{\pred{}}_{\writ}$~\eqref{eq:op-mem-comm-write}, and on any other step
it is left uninterpreted.  The memory vectors are then reconstructed
by the rule~\eqref{eq:mem-reconstruct}: $\mem_{k+1}:=\mem_k[\addr\mapsto x]$
with $\addr:=\hat{S}_k.\regs[0]$, $x:=\hat{S}_k.\regs[1]$ on a write step, and
$\mem_{k+1}:=\mem_k$ otherwise.

We proceed by induction on $k$, carrying the commitment invariant
$\hat{S}_k.\mem=\mathsf{Commit}(\mem_k)$.  It holds at $k=0$ as noted above.
For the inductive step, we verify that all requirements of
Proposition~\ref{prop:memory-extractability} are satisfied at step $k$:
\begin{enumerate}
  \item \textbf{Commitment invariant at $k$.}
    $\hat{S}_k.\mem=\mathsf{Commit}(\mem_k)$ is the induction hypothesis.
    Note that the proposition requires the invariant only at the \emph{first}
    of the two states; the invariant at $k+1$ is one of its conclusions, and is
    what carries the induction forward.
  \item \textbf{Step predicate.}
    From Step~5,
\[
  \widehat{\varphi}_{\step}(\hat{S}_k,\hat{S}_{k+1},\widehat{\BBB}_{\lceil(k+1)/N_{\seg}\rceil})=\mathrm{TRUE}.
\]
  \item \textbf{Opening witnesses (read case).}
    If step $k$ is a read, the witness $\pi$ parsed from $\pi^{\mathsf{mem}}_k$
    satisfies $\hat{S}_k.\mem=\hat{S}_{k+1}.\mem$ and
\[
  \mathsf{Verify}(\hat{S}_k.\mem,\hat{S}_k.\regs[0],\hat{S}_{k+1}.\regs[1],\pi)=1,
\]
    as required by the read requirement of the proposition.
  \item \textbf{Opening witnesses (write case).}
    If step $k$ is a write, the witnesses $(\pi_{\writ},v_{\mathsf{old}})$
    parsed from $\pi^{\mathsf{mem}}_k$ satisfy
\[
  \mathsf{Verify}(\hat{S}_k.\mem,\hat{S}_k.\regs[0],v_{\mathsf{old}},\pi_{\writ})=1
\]
and
\[
  \mathsf{Verify}(\hat{S}_{k+1}.\mem,\hat{S}_k.\regs[0],\hat{S}_k.\regs[1],\pi_{\writ})=1,
\]
    as required by the write requirement of the proposition.
\end{enumerate}
Applying Proposition~\ref{prop:memory-extractability} with the adversary
$\mathcal{A}_k$ defined as the composed extractor from Steps~1--5 that outputs
\[
  (\hat{S}_k,\hat{S}_{k+1},\mem_k,\widehat{\BBB}_{\lceil(k+1)/N_{\seg}\rceil},w_k),
\]
where $w_k$ is the memory-opening witness $\pi^{\mathsf{mem}}_k$ parsed above,
yields both conclusions for $S_k:=(\hat{S}_k.\pc,\hat{S}_k.\regs,\mem_k)$:
\[
  \hat{S}_{k+1}.\mem=\mathsf{Commit}(\mem_{k+1})
  \qquad\text{and}\qquad
  \varphi_{\step}(S_k,S_{k+1})=\mathrm{TRUE}.
\]
The first re-establishes the invariant at $k+1$ and completes the induction;
the second is the per-step obligation of~\eqref{eq:relation-star}.  Both are
obtained from the same pair of binding events, so the invariant is not paid for
separately.
The proposition invokes position-binding (for the read and write cases)
and update-binding (for the write case) of $\Com_{\mathsf{mem}}$; if either
property is violated at step $k$ this sub-step fails.  There are $T$
such steps, each contributing at most
\[
  \mathsf{Adv}^{\mathsf{pos}}_{\Com_{\mathsf{mem}}}(\mathcal{D}^{\mathsf{pos}}_k)
+\mathsf{Adv}^{\mathsf{upd}}_{\Com_{\mathsf{mem}}}(\mathcal{D}^{\mathsf{upd}}_k)
\]
to the failure probability.
Finally, the boundary memories match the claimed states: the extractor receives
$S_0$, so $\mem_0'=\mem_0$ by construction, and for the final memory the
reconstruction gives $\hat{S}_T.\mem=\mathsf{Commit}(\mem_T')$ while by hypothesis
$\hat{S}_T.\mem=\mathsf{Commit}(\mem_T)$; if $\mem_T'\neq\mem_T$, honest openings at
a differing index verify against the common root and break position-binding, so
$\mem_T'=\mem_T$ except with probability
\[
  \mathsf{Adv}^{\mathsf{pos}}_{\Com_{\mathsf{mem}}}(\mathcal{D}^{\mathsf{pos}}_{\mathsf{fin}}).
\]
This is a distinct bad event from the $T$ per-step ones --- it compares the
reconstructed $\mem_T'$ against the adversary's claimed $\mem_T$, which no
per-step reduction ever inspects --- and $\mathcal{D}^{\mathsf{pos}}_{\mathsf{fin}}$
is a correspondingly distinct reduction, which computes the two honest openings
at a differing index itself.  It therefore contributes its own summand
to~\eqref{eq:cte-bound} and is \emph{not} absorbed into the $\sum_k$ term.
At segment boundaries, $\hat{S}^{(i,N_{\seg})}=\hat{S}^{(i+1,0)}$ implies
$\mathsf{Commit}(\mem_{N_{\seg}\cdot i})$ is the same on both sides, so the
memory is consistent.

\textbf{Conclusion.}
Summing all failure contributions across Steps~1--6 gives exactly the
right-hand side of~\eqref{eq:cte-bound}:
\begin{align*}
  \Pr[\text{fail}]
  &\leq \mathsf{Adv}^{\mathsf{ks}}_{\Pi_4}(\mathcal{D}_4)
    &&\text{(Step 1)}\\
  &\quad+(m{-}1)\,\mathsf{Adv}^{\mathsf{ks}}_{\Pi_3}(\mathcal{D}_3)
    &&\text{(Step 2)}\\
  &\quad+m\,\mathsf{Adv}^{\mathsf{ks}}_{\Pi_2}(\mathcal{D}_2)
    &&\text{(Step 3)}\\
  &\quad+m\!\left(
      \mathsf{Adv}^{\mathsf{ks}}_{\Pi_1}(\mathcal{D}_1)
      +\sum_{j\in\{\step,\keccak,\poseidon,\range\}}\mathsf{Adv}^{\mathsf{ks}}_{\Pi_{0,j}}(\mathcal{D}_{0,j})
      +\mathsf{Adv}^{\mathsf{cr}}_{\Com_{\mathsf{bus}}}(\mathcal{D}_{\mathsf{bus}})
    \right)
    &&\text{(Step 4)}\\
  &\quad+\sum_{k=1}^{T}\!\left(
      \mathsf{Adv}^{\mathsf{pos}}_{\Com_{\mathsf{mem}}}(\mathcal{D}^{\mathsf{pos}}_k)
      +\mathsf{Adv}^{\mathsf{upd}}_{\Com_{\mathsf{mem}}}(\mathcal{D}^{\mathsf{upd}}_k)
    \right)
  +\mathsf{Adv}^{\mathsf{pos}}_{\Com_{\mathsf{mem}}}(\mathcal{D}^{\mathsf{pos}}_{\mathsf{fin}}).
    &&\text{(Step 6)}
\end{align*}
All coefficients $m = T/N_{\seg}$ and $T$ are polynomially bounded in $\lambda$,
so the right-hand side is a polynomially-sized sum of security advantages and is
negligible whenever each advantage term is.

\medskip
\noindent\textit{Running time.}
$\mathcal{E}_{\mathsf{zkVM}}$ never runs $\mathcal{A}$; its cost is the
straight-line sub-extractor calls plus deterministic work.  Steps~1--3 make
$1+(m-1)+m=2m$ extractor calls ($\mathcal{E}_{\Pi_4}$, one
$\mathcal{E}_{\Pi_3}$ per internal node, one $\mathcal{E}_{\Pi_2}$ per
segment); Step~4 makes $5m$ ($\mathcal{E}_{\Pi_1}$ and the four inner
extractors, per segment); Steps~5--6 are deterministic --- indexing and the
memory reconstruction rule~\eqref{eq:mem-reconstruct} --- costing
$\mathsf{poly}(\lambda)$ and invoking no extractor.  Under the convention that
each sub-extractor call costs at most
$\mathsf{Time}(\mathcal{A})+\mathsf{poly}(\lambda)$ (the ``Advantage
notation'' paragraph), $\mathsf{Time}(\mathcal{E}_{\mathsf{zkVM}}) \le
7m\cdot\mathsf{Time}(\mathcal{A})+\mathsf{poly}(\lambda)$.  Each reduction
adversary runs $\mathcal{A}$ once and re-runs the extraction chain above its
own layer, hence runs in at most
$(7m+1)\cdot\mathsf{Time}(\mathcal{A})+\mathsf{poly}(\lambda)$; this is the
constant $c=7m+1$ of the statement.  (Strictly, the convention of
Definition~\ref{def:extractable} bounds each extractor by the time of the
\emph{composed} adversary that produced its input, which compounds along the
recursion to $O(m^2)\cdot\mathsf{Time}(\mathcal{A})$ in the worst case; the
linear figure uses the per-call convention and is part of the idealization of
Remark~\ref{rem:idealized}.)
In particular,
$\mathcal{E}_{\mathsf{zkVM}}$ is PPT whenever $\mathcal{A}$ is PPT.
\end{proof}
